
% ======================================================================
\chapter{Introduzione}
\label{sec:intro}
% ======================================================================

Un linguaggio di programmazione è \textbf{reversibile} quando ogni sua
computazione può essere ripercorsa all'indietro fino a ritrovare esattamente lo
stato di partenza. Non si tratta di registrare una traccia da cui ricostruire il
passato: in un linguaggio reversibile l'informazione necessaria a tornare
indietro è già contenuta nello stato corrente, perché ogni costrutto è
progettato per non distruggerla. L'assegnamento distruttivo $x := e$, che
cancella il valore precedente di $x$, non appartiene a un linguaggio di questo
tipo; al suo posto compaiono aggiornamenti invertibili come $x \mathrel{+}= e$,
il cui inverso $x \mathrel{-}= e$ è immediato.

L'interesse per questo modo di calcolare ha tre radici indipendenti. La prima è
fisica: il principio di Landauer lega la cancellazione di un bit a una
dissipazione minima di energia, e un calcolo che non cancella può in linea di
principio evitarla. La seconda è pragmatica: molti problemi si presentano
naturalmente in coppie di trasformazioni inverse (compressione e
decompressione, codifica e decodifica, serializzazione e parsing) e un
linguaggio reversibile permette di scrivere \emph{un solo} programma invece di
due, eliminando per costruzione la possibilità che le due metà divergano. La
terza riguarda l'affidabilità: poter tornare indietro è il meccanismo di base
del \emph{recovery}, e infatti il \emph{rollback} di una transazione non è altro
che esecuzione all'indietro fino a un punto sicuro.

\medskip
\noindent
È proprio la terza radice a rendere interessante il caso \textbf{concorrente},
ed è qui che si colloca questo lavoro. In un sistema sequenziale l'esecuzione
all'indietro è una nozione semplice, perché c'è un solo passato da ripercorrere.
In presenza di più thread il passato non è più lineare: due rami che procedono
in parallelo possono aver interagito, e disfare un passo di uno può richiedere
di disfare prima i passi dell'altro che ne dipendono. Il capostipite dei
linguaggi reversibili, \textbf{Janus}~\cite{YokoyamaGluck2007}, è puramente
sequenziale; la reversibilità in contesto concorrente è stata studiata
soprattutto su calcoli di processi~\cite{LaneseMezzinaSchmittStefani2011,%
LaneseLienhardtMezzinaSchmittStefani2013}, che sono modelli e non linguaggi di
programmazione. Il territorio in mezzo (un linguaggio imperativo strutturato,
con la sintassi e i costrutti che un programmatore si aspetta, che sia insieme
reversibile e concorrente) è poco popolato: i tentativi esistono, ma finora
hanno definito i costrutti concorrenti attraverso la loro traduzione, senza una
semantica formale su cui poggiare le proprietà che ci si aspetta da un
linguaggio reversibile.

\medskip
\noindent
\textbf{Kairos} sta in quel territorio. Il nome è il greco
\kairosgr{} (\emph{kairós}): non il tempo che scorre, che è \chronosgr{}
(\emph{chrónos}), ma il momento opportuno, l'istante in cui una cosa va fatta. La distinzione è
pertinente a un linguaggio reversibile, dove il tempo non è una freccia ma una
coordinata percorribile nei due versi, e ciò che conta di un passo non è
\emph{quando} avviene ma che sia quello giusto: l'inversa di un programma non ne
riavvolge la storia, ne esegue un altro che riporta allo stato di partenza.

Al nucleo di Janus (assegnamenti invertibili, condizionali a doppia guardia,
cicli \texttt{from\,\dots\,until}, procedure invocabili con \texttt{call} e
\texttt{uncall}, interi e array) Kairos aggiunge il parallelismo esplicito
\texttt{par\,\dots\,and\,\dots\,rap}, la comunicazione sincrona su canali
tipizzati \texttt{ssend}/\texttt{srecv}, uno stack come struttura dati primitiva
e le transazioni \texttt{try\,\dots\,rollback\,\dots\,yrt}.
La scelta di progetto che tiene insieme il tutto è la \emph{disgiunzione}: rami
paralleli operano su porzioni disgiunte dello store e comunicano solo per
canali. Da questa singola ipotesi discende la confluenza degli interleaving che
non toccano canali (\S\ref{sec:par-conf}), e quindi il fatto che il
non-determinismo residuo non sia osservabile sullo store finale: la
concorrenza di Kairos serve a esprimere il parallelismo dei problemi, non a
introdurre corse critiche.

\medskip
\noindent
Il contributo di questo documento è la semantica formale del linguaggio.
Diamo una semantica operazionale strutturale in stile small-step, definita come
\textbf{estensione} della semantica di Janus di Lami, Lanese e
Stefani~\cite{LamiLaneseStefani2024}: ereditiamo da quella base la forma dei
giudizi, i contesti di valutazione, la configurazione terminale e il
trattamento degli errori, e aggiungiamo le regole per i costrutti che Kairos
introduce. Definiamo poi l'operatore di inversione $\Inv{\cdot}$, sintattico e
calcolabile in tempo lineare, e mostriamo in che senso realizza l'esecuzione
all'indietro.

\section{Notazione e convenzioni}
\label{sec:notazione}

I giudizi hanno la forma $\cfg{c}{\sigma} \to \cfg{c'}{\sigma'}$, con lo store a
sinistra come nella semantica di base, non il turnstile big-step
$\sigma \vdash e \Downarrow v$ usuale nella letteratura su Janus. Non è una
variante stilistica: interleaving dei rami e sincronizzazione sui canali
(\cref{sec:par}) si esprimono molto più naturalmente passo per passo, perché un
big-step collassa l'intera esecuzione di un comando in un unico giudizio e i
punti in cui intercalare l'altro ramo restano impliciti. È la stessa motivazione
dichiarata in~\cite{LamiLaneseStefani2024}, che adotta lo small-step per Janus
proprio in vista di una futura estensione concorrente.

Una convenzione tipografica riguarda ciò che sta \emph{fuori} dalla relazione di
transizione. Le regole d'errore, che producono $\err$, portano il nome in
\textcolor{red!60!black}{rosso}. I vincoli di buona formazione sono distinti
secondo il momento in cui li si verifica: \static{in verde} quelli controllati
sul testo del programma, prima di eseguirlo; \dynamic{in arancio} quelli
sorvegliati durante l'esecuzione, perché riguardano stato che il testo del
programma non determina.

\section{Struttura della tesi}

Il \cref{chap:background} richiama Janus, la semantica small-step su cui ci
basiamo e i lavori sulla reversibilità concorrente. Il \cref{chap:semantica} è
il nucleo tecnico: sintassi astratta, sistema di transizione, il confine fra ciò
che è ereditato dalla base e ciò che Kairos aggiunge, e le regole di riduzione
dagli assegnamenti fino a \texttt{par} e alle transazioni, con gli errori a
runtime. Il \cref{chap:reversibilita} definisce l'operatore di inversione, ne
dimostra le proprietà e riassume in un quadro unico tutte le regole. Seguono il
\cref{chap:programmare} sull'uso del linguaggio, il \cref{chap:interprete}
sull'implementazione, il \cref{chap:correlati} sul confronto con i lavori
esistenti e il \cref{chap:conclusioni} con le conclusioni.

% ======================================================================
\chapter{Background}
\label{chap:background}
\label{sec:back}
% ======================================================================

\subsubsection{Janus e la programmazione reversibile}

\textbf{Janus} è stato definito nel 1982 da Lutz e Derby come esercizio
didattico e riportato all'attenzione della comunità da Yokoyama e
Glück~\cite{YokoyamaGluck2007}, che ne hanno dato la prima semantica formale e
un self-interpreter invertibile. È un linguaggio imperativo strutturato in cui
ogni costrutto è invertibile per costruzione, secondo due idee ricorrenti.

La prima è l'\textbf{aggiornamento invertibile}. Al posto dell'assegnamento
distruttivo Janus ha solo aggiornamenti della forma $x \mathrel{\odot}= e$ con
$\odot \in \{+,-,\hat{\ }\}$, dove $x$ non compare in $e$: il vecchio valore di
$x$ è recuperabile dal nuovo applicando l'operazione inversa, e la condizione
sull'occorrenza di $x$ garantisce che il valore di $e$ sia lo stesso prima e
dopo il passo.

La seconda è la \textbf{doppia guardia}. Un condizionale irreversibile
$\mathtt{if}\ e\ \mathtt{then}\ c_1\ \mathtt{else}\ c_2$ non è invertibile
perché, arrivati alla fine, non si sa più quale ramo sia stato preso. Janus
richiede allora una seconda espressione, l'\emph{asserzione} in coda:
\[
\mathtt{if}\ e_1\ \mathtt{then}\ c_1\ \mathtt{else}\ c_2\ \mathtt{fi}\ e_2 ,
\]
con l'obbligo che $e_2$ valga vero all'uscita se e solo se si è preso il ramo
\texttt{then}. All'indietro è $e_2$ a decidere il ramo ed $e_1$ a fare da
asserzione: il ruolo delle due espressioni si scambia. Lo stesso schema governa
il ciclo
\[
\mathtt{from}\ e_1\ \mathtt{do}\ c_1\ \mathtt{loop}\ c_2\ \mathtt{until}\ e_2 ,
\]
dove $e_1$ deve valere solo alla prima iterazione ed $e_2$ solo all'ultima, così
che la condizione di ingresso all'indietro sia $e_2$ e quella di uscita $e_1$.
I due corpi sono eseguiti secondo la traccia $c_1\,[\,c_2\ c_1\,]^{*}$: $c_1$
porta avanti il lavoro dell'iterazione, $c_2$ fa avanzare il contatore fra
un'iterazione e la successiva. Kairos adotta la stessa forma; il ciclo a un
corpo solo ne è il caso particolare $c_2 = \mathtt{skip}$.

Su queste basi Janus è \emph{r-Turing completo}: calcola esattamente le funzioni
iniettive computabili, senza generare spazzatura. Le sue procedure si invocano
con \texttt{call} in avanti e \texttt{uncall} all'indietro, e l'inverso di un
programma si ottiene per trasformazione puramente sintattica.

\subsubsection{Una semantica small-step per Janus}

Le presentazioni classiche di Janus, a partire da~\cite{YokoyamaGluck2007},
usano una semantica \emph{big-step}: un giudizio $\sigma \vdash c \Downarrow
\sigma'$ mette in relazione lo store prima e dopo l'esecuzione \emph{completa}
di $c$. È una scelta comoda per un linguaggio sequenziale, ma collassa in un
unico passo l'intera computazione, e con essa i punti in cui un'altra
esecuzione potrebbe intercalarsi.

Lami, Lanese e Stefani~\cite{LamiLaneseStefani2024} hanno dato una semantica
\textbf{small-step} di Janus, con giudizi della forma $\langle \sigma, c\rangle
\to \langle \sigma', c'\rangle$ fra configurazioni, contesti di valutazione per
espressioni e comandi, configurazione terminale $\langle \sigma,
\mathtt{skip}\rangle$ ed errori segnalati da un esito distinto $\bot$. Uno degli
scopi dichiarati di quel lavoro è proprio predisporre il terreno per estensioni
concorrenti, sulle quali le proprietà di interesse si definiscono più
agevolmente in forma small-step. Kairos parte da lì: adottarne le convenzioni
significa che le scelte di presentazione sono già state fatte e argomentate a
monte, e che le regole nuove sono solo quelle dei costrutti nuovi. La
\S\ref{sec:base} elenca esattamente il confine fra le due cose.

Una semantica alternativa per Janus, in stile \emph{reversible}, è quella di
Lanese e Vidal~\cite{LaneseVidal2026}, in cui le regole in avanti e quelle
all'indietro sono date insieme sullo stesso sistema di transizione anziché
derivare le seconde da un inverter sintattico.

\subsubsection{Reversibilità e concorrenza}

Fuori dai linguaggi imperativi, la reversibilità in presenza di concorrenza è
stata studiata a fondo sui \emph{calcoli di processi}. Il punto delicato è che
in un sistema concorrente non esiste un unico ordine di esecuzione da
rovesciare: la nozione corretta è la \emph{reversibilità causalmente
consistente}, per cui un passo può essere disfatto solo dopo che sono stati
disfatti tutti i passi che da esso dipendono causalmente. In quest'ottica la
reversibilità va anche \emph{controllata}, perché tornare indietro senza limiti
non è utile a nulla: in $\rho\pi$ e in
croll-$\pi$~\cite{LaneseMezzinaSchmittStefani2011,%
LaneseLienhardtMezzinaSchmittStefani2013} il ritorno all'indietro è governato da
primitive esplicite, e il meccanismo che ne risulta (si torna indietro fino a
un punto scelto e da lì si prosegue diversamente) è esattamente quello delle
transazioni con compensazione. Le transazioni
\texttt{try\,\dots\,rollback\,\dots\,yrt} di Kairos (\S\ref{sec:try}) sono la
versione a linguaggio di programmazione della stessa idea.

Kairos adotta però una strada più restrittiva, e per questo più semplice: rami
paralleli lavorano su porzioni \emph{disgiunte} dello store e interagiscono
soltanto per comunicazione sincrona. Non c'è memoria condivisa, quindi non c'è
dipendenza causale da tracciare fra passi che non comunicano, e la consistenza
causale si riduce alla confluenza dimostrata in \S\ref{sec:par-conf}. Il prezzo
è che alcuni schemi di programmazione a memoria condivisa non sono esprimibili;
il guadagno è che nessun programma Kairos ben tipato può contenere una corsa
critica.

Fra i calcoli di processo e i linguaggi di programmazione esiste comunque un
precedente diretto: Janus è già stato esteso con un costrutto di composizione
parallela e uno scambio di messaggi, e l'estensione è stata realizzata con un
compilatore verso C++~\cite{Sammarchi2019}. Quel lavoro definisce i costrutti
attraverso il codice in cui vengono tradotti, non attraverso regole di
riduzione, e mantiene uno stato condiviso fra i rami; è quindi il punto di
partenza più vicino a quello di questa tesi, e il \cref{chap:correlati} lo
riprende per un confronto puntuale.

\subsubsection{Canali e tipi di sessione}

La comunicazione di Kairos avviene su canali \emph{sincroni} e \emph{tipizzati},
con il vincolo che ogni canale abbia in ogni istante al più due titolari, una
\emph{sessione binaria}. È una restrizione che viene dalla teoria dei
\textbf{tipi di sessione}~\cite{HondaVasconcelosKubo1998,%
HondaYoshidaCarbone2008}, dove il tipo di un canale non descrive solo i dati
scambiati ma la \emph{sequenza} degli scambi, e i tipi dei due estremi devono
essere duali. In Kairos il vincolo di sessione binaria serve a garantire che il
partner di ogni comunicazione sia determinato: senza
di esso un $\mathtt{ssend}$ potrebbe sincronizzarsi con più $\mathtt{srecv}$
diversi, e la scelta sarebbe un non-determinismo osservabile sullo store, in
contrasto con la confluenza.

Il vincolo ha due forme, e la scelta fra le due non è di dettaglio. Nella forma
\emph{statica} un canale appartiene agli stessi due rami per tutta la sua vita,
e questo si legge dal testo del programma. Nella forma \emph{dinamica} un canale
può stare prima fra due rami e poi fra altri due, purché il primo titolare
smetta di usarlo prima che il nuovo cominci: è ciò che nella teoria dei tipi di
sessione si chiama \textbf{delega}, il meccanismo con cui, in stile
$\pi$-calcolo, un processo passa a un terzo il canale su cui sta parlando. La
seconda è strettamente più espressiva, e Kairos adotta quella: il vincolo è
sorvegliato osservando l'esecuzione (\S\ref{sec:chan}), non leggendo il testo del
programma, perché il testo non basta a seguire un canale che cambia titolare.

\medskip
\noindent
Va detto dove sta il confine di questo lavoro. Il vincolo di sessione binaria e
l'ipotesi di disgiunzione sono qui \emph{assunzioni} sotto cui vale la
confluenza, non un sistema di tipi: la semantica delle
\S\S\ref{sec:esp}--\ref{sec:cmd} non impedisce di scrivere un programma con una
corsa critica, si limita a non garantire nulla su di esso. La separazione è
deliberata e segue la pratica corrente, in cui il calcolo si definisce senza
preoccuparsi delle corse e un sistema di tipi sovrapposto le esclude; i tipi di
sessione, binari o multiparty, sono precisamente una famiglia di sistemi di tipi
di questo genere. In \S\ref{sec:chan} diamo la disciplina binaria per i canali di
Kairos, delega compresa, e i tipi di sessione che ne descrivono il protocollo:
i tipi non si dichiarano ma si ricostruiscono dal programma, e i due estremi di
ogni canale devono risultare duali. Restano fuori la scelta etichettata e i
protocolli a più di due partecipanti, estensioni possibili che qui non
perseguiamo.

% ======================================================================
\chapter{Kairos: sintassi e semantica operazionale}
\label{chap:semantica}

\section{Sintassi astratta}
\label{sec:sintassi}
% ======================================================================

Insiemi di base: numeri interi $m \in \mathbb{Z}$, variabili $v \in \mathit{Var}$.
Categorie derivate:

\paragraph{Espressioni aritmetiche $e \in \mathit{Exp}$.}
\begin{align*}
e ::= m \mid v \mid a[e] \mid e \oplus e
\qquad
\oplus \in \{+,\,-,\,*,\,/,\,\%\}
\end{align*}

\noindent
$a[e]$ è la lettura della cella di indice $e$ dell'array $a$. Gli operatori
moltiplicativi compaiono solo dentro le espressioni, mai come assegnamento:
$*$ e $/$ non hanno un'inversa che restituisca l'operando, mentre a destra di
$\mathbin{+\!\!=}$ non tolgono nulla, perché l'inversa dell'assegnamento è
$\mathbin{-\!\!=}$ con la \emph{stessa} espressione, e questa dipende da
variabili che il comando non tocca. $/$ e $\%$ sono la divisione intera e il
resto; il divisore nullo non ha risultato e produce un errore:
\[
\dfrac{}{\cfg{m / 0}{\sigma} \toe \err}\ \rnameerr{Div-Err}
\]
con la regola gemella per $\%$. È un fallimento come gli altri
(\S\ref{sec:err}), non un valore convenzionale: un valore convenzionale
renderebbe l'operazione non iniettiva proprio dove serve che lo sia.

\paragraph{Condizioni $b \in \mathit{Cond}$.}
\begin{align*}
b ::= e = e \mid e \neq e \mid e \geq e \mid e \leq e \mid e > e \mid e < e
\end{align*}

\paragraph{Comandi $c \in \mathit{Com}$.}
\begin{align*}
c ::={}& \mathtt{skip} && \text{(comando vuoto)}\\
   \mid{}& \ell \mathbin{+\!\!=} e \mid \ell \mathbin{-\!\!=} e \mid \ell \mathbin{\hat{}\!=} e \mid \ell \mathbin{<\!=\!>} \ell && \text{(assegnamenti reversibili)}\\
   \mid{}& \mathtt{local}\ v = m \mid \mathtt{delocal}\ v = m && \text{(scope)}\\
   \mid{}& \mathtt{local}\ a[n] = m \mid \mathtt{delocal}\ a[n] = m && \text{(scope, array)}\\
   \mid{}& c\ ;\ c && \text{(sequenza)}\\
   \mid{}& \mathtt{if}\ b\ [\mathtt{then}\ c]\ [\mathtt{else}\ c]\ \mathtt{fi}\ b' && \text{(condizionale)}\\
   \mid{}& \begin{array}[t]{@{}l@{}}\mathtt{from}\ b_1\ [\mathtt{do}\ c_1]\\ \quad[\mathtt{loop}\ c_2]\ \mathtt{until}\ b_2\end{array} && \text{(ciclo)}\\
   \mid{}& \mathtt{call}\ f(\bar v) \mid \mathtt{uncall}\ f(\bar v) && \text{(chiamata / inversa)}\\
   \mid{}& \mathtt{push}(v,s) \mid \mathtt{pop}(v,s) && \text{(stack)}\\
   \mid{}& \mathtt{ssend}(\langle \bar v \rangle, \mathit{ch}) \mid \mathtt{srecv}(\langle \bar v \rangle, \mathit{ch}) && \text{(canali)}\\
   \mid{}& \mathtt{par}\ c\ \mathtt{and}\ c\ [\mathtt{and}\ c]^{*}\ \mathtt{rap} && \text{(parallelo)}\\
   \mid{}& \mathtt{try}\ b\ c\ [\mathtt{rollback}\ c]\ \mathtt{yrt}\ b' && \text{(commit / rollback)}
\end{align*}

\noindent
Il bersaglio di un assegnamento è un \emph{luogo} $\ell ::= v \mid a[e]$: una
variabile intera oppure una cella di array. Le due forme si comportano allo
stesso modo — l'unica differenza è che la cella va individuata valutando
l'indice prima di scrivere.

\noindent
Un \emph{programma} è una lista di procedure
$\mathtt{procedure}\ f(\bar v)\ \{\, c \,\}$; $\bar v$ è la lista (eventualmente
vuota) dei parametri. La sequenza è scritta con il punto e virgola esplicito,
$c_1\,;\,c_2$, come nelle regole di \S\ref{sec:seq}. Il condizionale ha
\textbf{due} guardie ($b$ d'ingresso, $b'$ d'uscita) e il ciclo pure ($b_1$
d'ingresso, $b_2$ d'uscita): è ciò che rende il linguaggio reversibile
(\S\ref{sec:rev}). Le parentesi quadre segnano le clausole \emph{facoltative},
come in Janus: un ramo assente vale $\mathtt{skip}$, e nelle regole che seguono
si scrivono per esteso tutte e due, senza perdita di generalità. Un programma ben formato soddisfa inoltre due vincoli: su
$\mathtt{par}$, ogni canale usato al suo interno è condiviso da esattamente
due rami concorrenti (\S\ref{sec:chan}); su $\mathtt{try}\ldots
\mathtt{rollback}\ldots\mathtt{yrt}$, né il corpo $c$ né il corpo di
rollback $r$ possono contenere $\mathtt{ssend}$/$\mathtt{srecv}$, né
direttamente né tramite $\mathtt{call}$/$\mathtt{uncall}$
(\S\ref{sec:try}).

% ======================================================================
\section{Sistema di transizione e store}
\label{sec:store}
% ======================================================================

Un sistema di transizione è una tripla $\langle \Gamma, \top, \to \rangle$ con
$\Gamma$ insieme degli stati, $\top \subseteq \Gamma$ stati terminali,
${\to} \subseteq \Gamma \times \Gamma$ relazione di transizione. $\to^*$ è la
sua chiusura riflessiva e transitiva:
\[
\dfrac{}{y \to^* y}
\qquad
\dfrac{y \to^* y' \quad y' \to y''}{y \to^* y''}
\]

\noindent
Per dare significato alle variabili introduciamo uno \textbf{store}. In Kairos
le variabili sono di quattro tipi: \texttt{int} (interi), \texttt{int}$[n]$
(array: $n$ interi, con $n$ fissato alla dichiarazione), \texttt{stack}
(sequenze di interi) e \texttt{channel} (canali: nomi di sincronizzazione tra
rami paralleli). Array e stack sono entrambi sequenze di interi ma si usano in
modi opposti: l'array ha lunghezza fissa e si legge per indice senza
consumarlo, lo stack cresce e cala e si tocca solo in cima. Un canale \textbf{non} contiene valori: è un punto di
rendez-vous, non un contenitore (\S\ref{sec:chan}), a differenza di
\texttt{stack}, un canale allocato mantiene sempre la stessa identità per
tutta la sua vita, e $\sigma(\mathit{ch})$ non compare mai
nell'aggiornamento di una regola. Indicando con $\mathit{Chan}$ l'insieme
(opaco) delle identità di canale, lo store è
\[
\sigma : \mathit{Var} \rightharpoonup \mathbb{Z} \cup \mathbb{Z}^{n} \cup \mathbb{Z}^{*} \cup \mathit{Chan},
\]
parziale perché lo scope cambia con \texttt{local}/\texttt{delocal}. Per le
sequenze: $\varepsilon$ è la vuota, $n : \ell$ è il \emph{cons} (la testa $n$
è la cima dello stack). Scriviamo $\upd{\sigma}{v}{w}$ per l'aggiornamento e
$\sigma \setminus v$ per la rimozione di $v$ dal dominio. $\varepsilon$ denota
qui \textbf{solo} la sequenza vuota di uno \texttt{stack}: è un valore, non un
comando. Il comando vuoto è un simbolo distinto, $\mathtt{skip}$
(\S\ref{sec:sintassi}, \S\ref{sec:seq}), le due nozioni non vanno confuse.

% ======================================================================
\section{La semantica di base}
\label{sec:base}
% ======================================================================

Kairos non viene definito da zero. La sua semantica è data come
\textbf{estensione} della semantica operazionale small-step di Janus di Lami,
Lanese e Stefani~\cite{LamiLaneseStefani2024}: da lì ereditiamo la forma del
giudizio, la configurazione terminale, il trattamento degli errori e
l'apparato dei contesti di valutazione; a quella base aggiungiamo i costrutti
che Janus non ha. Le scelte di progetto già compiute in quel lavoro non
vengono quindi ridiscusse qui: dove questo documento tace, vale la convenzione
della base. Nel seguito chiamiamo \emph{la base} quella semantica, e
segnaliamo esplicitamente ogni punto in cui ce ne discostiamo.

\paragraph{Giudizio.} Le regole hanno la forma
\[
\cfg{c}{\sigma} \to \cfg{c'}{\sigma'}
\qquad\text{oppure}\qquad
\cfg{c}{\sigma} \to \bot ,
\]
dove la seconda denota un \textbf{errore a runtime} dovuto al fallimento di
un'asserzione. Lo store sta a \emph{sinistra} del comando, come nella base.
La base usa un'unica relazione $\to$ su tutte le categorie sintattiche; qui
manteniamo la decorazione $\toe$ / $\tob$ / $\toc$ per distinguere il
frammento delle espressioni, quello delle condizioni e quello dei comandi,
ma si tratta della stessa relazione letta su categorie diverse.

\paragraph{Configurazione terminale.} Nella base la computazione termina in
$\cfg{\mathtt{skip}}{\sigma}$, e per questo \emph{non} esiste una regola per
$\mathtt{skip}$: citando~\cite{LamiLaneseStefani2024}, \emph{``there is no
rule for \textup{\texttt{skip}}, since $\langle \sigma, \mathtt{skip}\rangle$
denotes the end of the computation''}. È la convenzione standard, e la
adottiamo: ogni regola che porta a termine un comando conclude in
$\cfg{\mathtt{skip}}{\sigma}$.

\medskip
\noindent
\textbf{Configurazione terminale.} La convenzione è adottata
\emph{ovunque} in questo documento: ogni regola che porta a termine un comando
conclude in $\cfg{\mathtt{skip}}{\sigma}$, non nello store nudo. Non esiste
quindi una regola per $\mathtt{skip}$; il consumo di $\mathtt{skip}$ in
sequenza è affidato a \textsc{SeqS} (\S\ref{sec:seq}), e la terminazione di un
ramo parallelo a \textsc{Par-L0}/\textsc{Par-R0} (\S\ref{sec:par}).

\paragraph{Errori.} L'esito d'errore è $\bot$, un elemento distinto da ogni
configurazione. Un errore non viene propagato da una regola per ciascun
costrutto, ma da un'unica regola contestuale \textsc{CtxError} (sotto): è la
scelta della base, e sostituisce le regole di propagazione ad hoc.

\paragraph{Contesti di valutazione.} Per dare una definizione compatta della
semantica, la base solleva la riduzione di sotto-termini tramite
\emph{contesti di valutazione}: un contesto è un comando in cui un
sotto-termine è sostituito da un buco $\bullet$. Si distinguono i contesti
d'\textbf{espressione} $C_e[\bullet]$, il cui buco va riempito con
un'espressione o una condizione, e i contesti di \textbf{comando}
$C_s[\bullet]$, il cui buco va riempito con un comando; $C_e[e]$ e $C_s[c]$
denotano il comando ottenuto sostituendo $\bullet$ rispettivamente con $e$ e
con $c$. Adattati a Kairos, i contesti sono
{\footnotesize
\[
\begin{aligned}
C_e[\bullet] ::={}&
  v \mathbin{+\!\!=} \bullet \mid v \mathbin{-\!\!=} \bullet \mid v \mathbin{\hat{}\!=} \bullet
  && \text{(\S\ref{sec:assign})}\\
 \mid{}& \mathtt{if}\ \bullet\ \mathtt{then}\ c_1\ \mathtt{else}\ c_2\ \mathtt{fi}\ b_2
  && \text{(\S\ref{sec:if})}\\
 \mid{}& \mathtt{if}\ \ttt\ \mathtt{then}\ \mathtt{skip}\ \mathtt{else}\ c_2\ \mathtt{fi}\ \bullet
  \mid \mathtt{if}\ \ff\ \mathtt{then}\ c_1\ \mathtt{else}\ \mathtt{skip}\ \mathtt{fi}\ \bullet & \\
 \mid{}& ((\bullet,b_1),\, c_1,\, b_2,\, c_2)
  \mid (b_1,\, (\mathtt{skip},c_1),\, (\bullet,b_2),\, c_2)
  && \text{(\S\ref{sec:loop})}\\
 \mid{}& ((\bullet,b_1),\, (\mathtt{skip},c_1),\, b_2,\, (\mathtt{skip},c_2)) & \\[0.4em]
C_s[\bullet] ::={}& \bullet\,;\,c
  && \text{(\S\ref{sec:seq})}\\
 \mid{}& \mathtt{if}\ \ttt\ \mathtt{then}\ \bullet\ \mathtt{else}\ c_2\ \mathtt{fi}\ b_2
  \mid \mathtt{if}\ \ff\ \mathtt{then}\ c_1\ \mathtt{else}\ \bullet\ \mathtt{fi}\ b_2
  && \text{(\S\ref{sec:if})}\\
 \mid{}& (b_1,\, (\bullet,c_1),\, (b_2,b_2),\, c_2)
  \mid ((b_1,b_1),\, (\mathtt{skip},c_1),\, b_2,\, (\bullet,c_2))
  && \text{(\S\ref{sec:loop})}\\
 \mid{}& \mathtt{par}\ \bullet\ \mathtt{and}\ c\ \mathtt{rap}
  \mid \mathtt{par}\ c\ \mathtt{and}\ \bullet\ \mathtt{rap}
  && \text{(\S\ref{sec:par}, propria di Kairos)}
\end{aligned}
\]
}
Le clausole con le quadruple appartengono al ciclo e sono spiegate in
\S\ref{sec:loop}; quelle di $\mathtt{par}$ sono l'unica aggiunta di Kairos,
e sono discusse in \S\ref{sec:par}. Non c'è una clausola per
$\mathtt{try}$: le sue regole (\S\ref{sec:try}) riducono il corpo con
$\toc^{*}$ in una premessa, non in posizione, quindi il corpo non è un buco
di contesto. Le tre regole contestuali sono quelle della base:
\[
\dfrac{\cfg{e}{\sigma} \toe \cfg{e'}{\sigma}}
      {\cfg{C_e[e]}{\sigma} \toc \cfg{C_e[e']}{\sigma}} \rname{CtxExp}
\qquad
\dfrac{\cfg{c}{\sigma} \toc \cfg{c'}{\sigma'}}
      {\cfg{C_s[c]}{\sigma} \toc \cfg{C_s[c']}{\sigma'}} \rname{CtxStmt}
\]
\[
\dfrac{\cfg{c}{\sigma} \toc \bot}
      {\cfg{C_s[c]}{\sigma} \toc \bot} \rnameerr{CtxError}
\]
\textsc{CtxExp} e \textsc{CtxStmt} sollevano un passo riuscito,
\textsc{CtxError} solleva un errore. Le stesse regole valgono per una
condizione al posto di un'espressione, sostituendo $\toe$ con $\tob$: nella
base gli operatori relazionali sono operatori binari come gli altri e non
esiste una categoria separata (\S\ref{sec:esp}). Applicandole ripetutamente si
attraversa un annidamento di contesti di profondità qualunque.

\medskip
\noindent
Queste tre regole rendono superflue tutte le regole di congruenza e di
propagazione dell'errore scritte una per costrutto: \textsc{Seq1}
(\S\ref{sec:seq}) è \textsc{CtxStmt} istanziata su $C_s[\bullet] =
\bullet\,;\,c$, \textsc{Par-L}/\textsc{Par-R} (\S\ref{sec:par}) sono
\textsc{CtxStmt} istanziata sulle due posizioni del $\mathtt{par}$, e
\textsc{CtxError} rende superflua ogni regola di
propagazione dell'errore dedicata a un singolo costrutto (\S\ref{sec:err}).
Le istanze di contesto le manteniamo scritte per esteso dove aiutano la
lettura, segnalando di volta in volta che sono istanze e non regole
indipendenti.

% ----------------------------------------------------------------------
\subsubsection*{Che cosa aggiunge Kairos}
% ----------------------------------------------------------------------

Quanto segue è il \textbf{contributo di questa tesi}: sono i costrutti e i
vincoli che la base non ha e che vengono definiti qui per la prima volta sopra
di essa. La base è esplicita sui propri limiti:
\emph{``In this paper, we do not consider parameters or local
variables''}~\cite{LamiLaneseStefani2024}. Kairos estende dunque quella
semantica su sette fronti:

\begin{enumerate}
\item \textbf{Variabili locali}: $\mathtt{local}$ / $\mathtt{delocal}$
  (\S\ref{sec:scope}), assenti dalla base.
\item \textbf{Parametri di procedura}: $\mathtt{call}\ f(\bar v)$ /
  $\mathtt{uncall}\ f(\bar v)$ (\S\ref{sec:call}); la base ha solo
  $\mathtt{call}\ \mathit{id}$ senza parametri, e passa i valori per effetto
  collaterale sullo store globale.
\item \textbf{Stack}: $\mathtt{push}$ / $\mathtt{pop}$ con la convenzione
  \emph{zero-cleared} (\S\ref{sec:stack}); la base ha array $x[e]$, non stack.
\item \textbf{Canali}: $\mathtt{ssend}$ / $\mathtt{srecv}$ con
  \textbf{transizioni etichettate} (\S\ref{sec:chan}).
\item \textbf{Composizione parallela}: $\mathtt{par}\ldots\mathtt{and}\ldots
  \mathtt{rap}$, con interleaving e sincronizzazione \textsc{Par-Comm}
  (\S\ref{sec:par}).
\item \textbf{Transazioni}: $\mathtt{try}\ldots\mathtt{rollback}\ldots
  \mathtt{yrt}$ (\S\ref{sec:try}).
\item \textbf{Sessione binaria} sui canali come proprietà dinamica, delega
  compresa (\S\ref{sec:chan}), con le proprietà che ne discendono: confluenza
  (\S\ref{sec:par}) e reversibilità estesa al caso $\mathtt{par}$
  (\S\ref{sec:rev}).
\end{enumerate}

\noindent
La base è sequenziale, ma adotta lo small-step \emph{proprio in vista} di
un'estensione concorrente: i punti 4--7 sono l'estensione che quella scelta
rendeva possibile.

% ======================================================================
\section{Espressioni e condizioni}
\label{sec:esp}
% ======================================================================

Adottiamo \textbf{senza modifiche} la semantica small-step delle espressioni
della base (\cite{LamiLaneseStefani2024}, Fig.~5): il giudizio
$\cfg{e}{\sigma} \toe \cfg{e'}{\sigma}$ non tocca mai lo store e riduce
costanti (\textsc{ConS}), variabili (\textsc{VarS}) e operatori binari
(\textsc{Bop1}--\textsc{Bop3}) con strategia interna sinistra. Come nella
base, la sintassi di espressioni e condizioni è \emph{estesa con i valori come
foglie}, così che un termine parzialmente valutato resti un termine: è ciò
che permette ai contesti $C_e[\bullet]$ (\S\ref{sec:base}) di sollevarne la
riduzione.

\medskip
\noindent
Le \textbf{condizioni} non sono una categoria separata nella base: gli
operatori relazionali ($=,\neq,\geq,\leq,>,<$) sono già fra i suoi operatori
binari $\odot$, e non richiedono quindi regole proprie. Manteniamo per
leggibilità la decorazione $\tob$ per la riduzione di una condizione, i cui
valori sono $\ttt$ (vero) e $\ff$ (falso), dove la base usa i predicati
$\mathsf{is\_true?}(v)$ e $\mathsf{is\_false?}(v)$ su un valore qualunque.

\medskip
\noindent
Poiché $\toe$ e $\tob$ sono deterministiche e non modificano lo store, si ha
$\cfg{e}{\sigma} \toe^{*} \cfg{m}{\sigma}$ con $m$ unico e
$\cfg{b}{\sigma} \tob^{*} \cfg{\beta}{\sigma}$ con $\beta \in \{\ttt,\ff\}$
unico; scriviamo $\mathit{eval}(e,\sigma) = m$ e
$\mathit{eval}_b(b,\sigma) = \beta$. In Kairos $\mathtt{push}$ e
$\mathtt{pop}$ restano \textbf{fuori} dalle espressioni
(\S\ref{sec:stack}): le espressioni non hanno effetti collaterali, e la loro
valutazione resta perciò esattamente quella della base.

% ======================================================================
\section{Semantica dei comandi ($\toc$)}
\label{sec:cmd}
% ======================================================================

$\Gamma_c = \{\cfg{c}{\sigma}\} \cup \{\bot\}$, terminali
$\top_c = \{\cfg{\mathtt{skip}}{\sigma}\} \cup \{\bot\}$: come nella semantica
di base (\S\ref{sec:base}), una configurazione \emph{terminale} è
$\cfg{\mathtt{skip}}{\sigma}$: non esiste alcuna regola per $\mathtt{skip}$,
perché $\cfg{\mathtt{skip}}{\sigma}$ \emph{è} la fine della computazione.

\subsubsection{Sequenza}
\label{sec:seq}
$\mathtt{skip}$ è il comando vuoto e, da solo, è la configurazione
terminale: non ha regole. Non va confuso con $\varepsilon$
(\S\ref{sec:store}), che è la sequenza vuota di uno \texttt{stack}, un
\emph{valore}, non un comando. In sequenza, $\mathtt{skip}$ viene consumato
dalla regola \textsc{SeqS} della base:
\[
\dfrac{\cfg{c_1}{\sigma} \toc \cfg{c_1'}{\sigma'}}
      {\cfg{c_1\,;\,c_2}{\sigma} \toc \cfg{c_1'\,;\,c_2}{\sigma'}} \rname{Seq1}
\qquad
\dfrac{}{\cfg{\mathtt{skip}\,;\,c_2}{\sigma} \toc \cfg{c_2}{\sigma}} \rname{SeqS}
\]
\[
\dfrac{\cfg{c_1}{\sigma} \xrightarrow{\ \alpha\ } \cfg{c_1'}{\sigma}}
      {\cfg{c_1\,;\,c_2}{\sigma} \xrightarrow{\ \alpha\ } \cfg{c_1'\,;\,c_2}{\sigma}} \rname{Seq-Lbl}
\]
\textsc{Seq1} è \textsc{CtxStmt} (\S\ref{sec:base}) istanziata sul contesto
$C_s[\bullet] = \bullet\,;\,c$: la riportiamo per esteso solo per leggibilità.
\textsc{SeqS} non è invece un'istanza di contesto: è la regola della base che
consuma un $\mathtt{skip}$ già raggiunto e fa avanzare la sequenza al comando
successivo (\S\ref{sec:base}). \textsc{Seq-Lbl}
solleva un passo \textbf{etichettato} (\S\ref{sec:chan}) di $c_1$ nella
sequenza, per un'etichetta $\alpha$ generica: serve perché \textsc{CtxStmt}
solleva solo passi $\toc$ non etichettati, ed è ciò che permette a un
$\mathtt{ssend}/\mathtt{srecv}$ seguito da altro codice (l'idioma normale,
non solo un canale da solo nel corpo di una procedura), di offrirsi comunque
come partner in \textsc{Par-Comm} (\S\ref{sec:par}). Non serve invece alcuna
regola d'errore per la sequenza: un $\err$ prodotto da $c_1$ risale grazie a
\textcolor{red!60!black}{\textsc{CtxError}} sullo stesso contesto
$\bullet\,;\,c$ (\S\ref{sec:err}).

\subsubsection{Assegnamenti reversibili}
\label{sec:assign}
\static{Vincolo (verificato staticamente): la variabile aggiornata non compare
a destra ($v \notin e$)}, così l'operazione è iniettiva.
\[
\dfrac{\cfg{e}{\sigma} \toe \cfg{e'}{\sigma}}
      {\cfg{v \mathbin{+\!\!=} e}{\sigma} \toc \cfg{v \mathbin{+\!\!=} e'}{\sigma}} \rname{Add1}
\]
\[
\dfrac{}{\cfg{v \mathbin{+\!\!=} m}{\sigma} \toc \cfg{\mathtt{skip}}{\upd{\sigma}{v}{\sigma(v)+m}}} \rname{Add2}
\]
\[
\dfrac{\cfg{e}{\sigma} \toe \cfg{e'}{\sigma}}
      {\cfg{v \mathbin{-\!\!=} e}{\sigma} \toc \cfg{v \mathbin{-\!\!=} e'}{\sigma}} \rname{Sub1}
\]
\[
\dfrac{}{\cfg{v \mathbin{-\!\!=} m}{\sigma} \toc \cfg{\mathtt{skip}}{\upd{\sigma}{v}{\sigma(v)-m}}} \rname{Sub2}
\]
\[
\dfrac{\cfg{e}{\sigma} \toe \cfg{e'}{\sigma}}
      {\cfg{v \mathbin{\hat{}\!=} e}{\sigma} \toc \cfg{v \mathbin{\hat{}\!=} e'}{\sigma}} \rname{Xor1}
\]
\[
\dfrac{}{\cfg{v \mathbin{\hat{}\!=} m}{\sigma} \toc \cfg{\mathtt{skip}}{\upd{\sigma}{v}{\sigma(v)\xor m}}} \rname{Xor2}
\]
\[
\dfrac{}{\cfg{v_1 \mathbin{<\!=\!>} v_2}{\sigma} \toc \cfg{\mathtt{skip}}{\upd{\upd{\sigma}{v_1}{\sigma(v_2)}}{v_2}{\sigma(v_1)}}} \rname{Swap}
\]

\paragraph{Celle di array.} Quando il bersaglio è $a[e_i]$ le regole sono le
stesse, con un passo in più davanti: l'indice si riduce a un intero prima che
si scriva. Scriviamo $\sigma(a)_k$ per la $k$-esima componente e
$\updc{\sigma}{a}{k}{m}$ per l'aggiornamento della sola cella $k$, cioè
$\upd{\sigma}{a}{\sigma(a)[k \mapsto m]}$. Diamo le regole per
$\mathbin{+\!\!=}$; $\mathbin{-\!\!=}$ e $\mathbin{\hat{}\!=}$ sono identiche a
meno dell'operatore.
\[
\dfrac{\cfg{e_i}{\sigma} \toe \cfg{e_i'}{\sigma}}
      {\cfg{a[e_i] \mathbin{+\!\!=} e}{\sigma} \toc \cfg{a[e_i'] \mathbin{+\!\!=} e}{\sigma}} \rname{IdxL}
\]
\[
\dfrac{0 \leq k < |\sigma(a)|}
      {\cfg{a[k] \mathbin{+\!\!=} m}{\sigma} \toc \cfg{\mathtt{skip}}{\updc{\sigma}{a}{k}{\sigma(a)_k + m}}} \rname{AddA}
\]
\[
\dfrac{k < 0 \ \lor\ k \geq |\sigma(a)|}
      {\cfg{a[k] \mathbin{+\!\!=} m}{\sigma} \toc \err} \rnameerr{Idx-Err}
\]

\noindent
L'indice fuori dai limiti è un errore e non un accesso silenzioso: in un
linguaggio reversibile scrivere fuori dal vettore romperebbe l'inversione senza
lasciare traccia, e l'errore comparirebbe altrove.

\noindent
Il vincolo statico si estende di conseguenza: in $a[e_i] \mathbin{+\!\!=} e$ il
nome $a$ non può comparire in $e$. È più forte del necessario — basterebbe che
non fosse la \emph{stessa} cella — ma gli indici non sono decidibili
staticamente, ed è la stessa scelta di Janus. Lo scambio $a[e_1]
\mathbin{<\!=\!>} a[e_2]$ resta invece ammesso: è involutivo comunque siano gli
indici, ed è quello che rende scrivibile un ordinamento sul posto.

\subsubsection{Scope: \texttt{local} / \texttt{delocal}}
\label{sec:scope}
\texttt{local}\ alloca $v$ con valore $m$; \texttt{delocal}\ \emph{verifica} che
$v$ valga ancora $m$ e poi lo rimuove (l'asserzione rende l'operazione
invertibile).
\[
\dfrac{v \notin \mathrm{dom}(\sigma)}
      {\cfg{\mathtt{local}\ v = m}{\sigma} \toc \cfg{\mathtt{skip}}{\upd{\sigma}{v}{m}}} \rname{Local}
\]
\[
\dfrac{\sigma(v) = m}
      {\cfg{\mathtt{delocal}\ v = m}{\sigma} \toc \cfg{\mathtt{skip}}{\sigma \setminus v}} \rname{Deloc}
\]
\[
\dfrac{v \in \mathrm{dom}(\sigma)}
      {\cfg{\mathtt{local}\ v = m}{\sigma} \toc \err} \rnameerr{Local-Err}
\]
\[
\dfrac{v \notin \mathrm{dom}(\sigma) \ \lor\ \sigma(v) \neq m}
      {\cfg{\mathtt{delocal}\ v = m}{\sigma} \toc \err} \rnameerr{Deloc-Err}
\]
Per un array il valore dichiarato vale per \emph{ogni} cella: aprirlo lo
riempie di $m$, chiuderlo verifica che tutte lo valgano ancora. Una sola cella
fuori posto è informazione non restituita, e la chiusura fallisce.
\[
\dfrac{a \notin \mathrm{dom}(\sigma)}
      {\cfg{\mathtt{local}\ a[n] = m}{\sigma} \toc \cfg{\mathtt{skip}}{\upd{\sigma}{a}{m^{n}}}} \rname{LocalA}
\]
\[
\dfrac{\sigma(a) = m^{n}}
      {\cfg{\mathtt{delocal}\ a[n] = m}{\sigma} \toc \cfg{\mathtt{skip}}{\sigma \setminus a}} \rname{DelocA}
\]
dove $m^{n}$ è la sequenza di $n$ copie di $m$. I due casi d'errore sono quelli
di \textsc{Local-Err} e \textsc{Deloc-Err}, con in più il disaccordo sulla
lunghezza: chiudere come $a[n']$ un array aperto come $a[n]$, con $n' \neq n$,
non è l'inversa dell'apertura.
Un'unica
sottigliezza propria di \textcolor{red!60!black}{\textsc{Deloc-Err}}: se $v \notin \mathrm{dom}(\sigma)$,
$\sigma(v)$ non è definito e non si può confrontare con $m$: la premessa
copre esplicitamente anche questo caso, non solo $\sigma(v) \neq m$.

\subsubsection{Condizionale (doppia guardia)}
\label{sec:if}
La guardia d'uscita $b'$ deve essere coerente col ramo: \emph{vera} dopo
\texttt{then}, \emph{falsa} dopo \texttt{else}. Adottiamo le quattro regole
della base (\cite{LamiLaneseStefani2024}, Fig.~8), che \textbf{non} usano
alcun comando ausiliario $\mathtt{assert}$: la guardia d'uscita si valuta
\emph{nel contesto}, in posizione, e il confronto fra il valore del test e
quello dell'asserzione decide fra passo riuscito ed errore.

\medskip
\noindent
Le regole si applicano quindi a un condizionale già ridotto: \textsc{CtxExp}
(\S\ref{sec:base}) ha portato il test a un valore, \textsc{CtxStmt} ha ridotto
a $\mathtt{skip}$ il ramo selezionato, e \textsc{CtxExp} ha infine valutato
l'asserzione. Le clausole di $C_e[\bullet]$ e $C_s[\bullet]$ per il
condizionale sono esattamente le posizioni usate qui.
\[
\dfrac{}{\cfg{\mathtt{if}\ \ttt\ \mathtt{then}\ \mathtt{skip}\ \mathtt{else}\ c_e\ \mathtt{fi}\ \ttt}{\sigma}
       \toc \cfg{\mathtt{skip}}{\sigma}} \rname{IfTrueEnd}
\]
\[
\dfrac{}{\cfg{\mathtt{if}\ \ff\ \mathtt{then}\ c_t\ \mathtt{else}\ \mathtt{skip}\ \mathtt{fi}\ \ff}{\sigma}
       \toc \cfg{\mathtt{skip}}{\sigma}} \rname{IfFalseEnd}
\]
\[
\dfrac{}{\cfg{\mathtt{if}\ \ttt\ \mathtt{then}\ \mathtt{skip}\ \mathtt{else}\ c_e\ \mathtt{fi}\ \ff}{\sigma}
       \toc \bot} \rnameerr{IfError1}
\]
\[
\dfrac{}{\cfg{\mathtt{if}\ \ff\ \mathtt{then}\ c_t\ \mathtt{else}\ \mathtt{skip}\ \mathtt{fi}\ \ttt}{\sigma}
       \toc \bot} \rnameerr{IfError2}
\]
Le prime due regole danno il passo riuscito quando test e asserzione hanno lo
\emph{stesso} valore di verità (entrambi $\ttt$, caso \textsc{IfTrueEnd};
entrambi $\ff$, caso \textsc{IfFalseEnd}); le altre due segnalano l'errore
quando i due valori \emph{differiscono}. L'asserzione d'uscita non
richiede quindi alcun comando ausiliario da eseguire in coda al ramo: viene
valutata in posizione di contesto, come il test d'ingresso.%
\footnote{Nella Fig.~8 di~\cite{LamiLaneseStefani2024} la regola
\textsc{IfFalseEnd} è stampata con premessa $\mathsf{is\_true?}(v_2)$, il che
la renderebbe identica a \textsc{IfError2}. Si tratta evidentemente di un
refuso: il testo che accompagna la figura dice esplicitamente che le prime due
regole coprono il caso in cui test e asserzione sono \emph{``either both true
\textup{(\textsc{IfTrueEnd}} case\textup{)} or both false
\textup{(\textsc{IfFalseEnd}} case\textup{)}''}. Qui riportiamo la versione
corretta, con l'asserzione falsa.}

\subsubsection{Ciclo \texttt{from\,\dots\,do\,\dots\,loop\,\dots\,until}}
\label{sec:loop}
Adottiamo il ciclo \textbf{a due corpi} della base, con le sue sette regole
(\cite{LamiLaneseStefani2024}, Fig.~9): $b_1$ è l'asserzione d'ingresso (vera
solo alla prima entrata), $c_1$ il \emph{do}-corpo, $c_2$ il \emph{loop}-corpo
eseguito \emph{fra} un'iterazione e la successiva, $b_2$ il test d'uscita. La
traccia d'esecuzione è $c_1\,[c_2\,c_1]^{*}$: $c_1$ gira $k$ volte e $c_2$
gira $k-1$ volte.

\medskip
\noindent
Come nella base, l'esecuzione usa un termine ausiliario a runtime, la
\textbf{quadrupla} $(X_1, X_2, X_3, X_4)$, dove ciascuna componente è o il
termine originale \emph{a riposo}, o una \emph{coppia}
$(\text{corrente},\allowbreak\ \text{originale})$ che segna la componente come in corso di valutazione e ne
conserva accanto l'originale, per poterlo riavviare all'iterazione seguente:
\[
\begin{array}{llll}
X_1 ::= b_1 \mid (\beta, b_1) &
X_2 ::= c_1 \mid (c, c_1) &
X_3 ::= b_2 \mid (\beta, b_2) &
X_4 ::= c_2 \mid (c, c_2)
\end{array}
\]
con $\beta$ un reduct di una condizione. La distinzione fra forma nuda e forma
a coppia è ciò che rende le regole \textbf{mutuamente esclusive}: p.es.\
\textsc{LoopExp1T} si applica quando $X_2$ e $X_4$ sono nudi (prima entrata),
\textcolor{red!60!black}{\textsc{LoopError2}} quando sono coppie con corrente
$\mathtt{skip}$ (fine di un'iterazione). Le riduzioni interne alla quadrupla
sono sollevate dai contesti $C_e[\bullet]$ e $C_s[\bullet]$ di
\S\ref{sec:base}, le cui clausole per il ciclo sono esattamente le posizioni
usate qui.
\[
\dfrac{}{\cfg{\mathtt{from}\ b_1\ \mathtt{do}\ c_1\ \mathtt{loop}\ c_2\ \mathtt{until}\ b_2}{\sigma}
       \toc \cfg{((b_1,b_1),\, c_1,\, b_2,\, c_2)}{\sigma}} \rname{LoopMainS}
\]
\[
\dfrac{}{\cfg{((\ttt,b_1),\, c_1,\, b_2,\, c_2)}{\sigma}
       \toc \cfg{(b_1,\, (c_1,c_1),\, (b_2,b_2),\, c_2)}{\sigma}} \rname{LoopExp1T}
\]
\[
\dfrac{}{\cfg{(b_1,\, (\mathtt{skip},c_1),\, (\ttt,b_2),\, c_2)}{\sigma}
       \toc \cfg{\mathtt{skip}}{\sigma}} \rname{LoopEndS}
\]
{\small
\[
\dfrac{}{\cfg{(b_1,\, (\mathtt{skip},c_1),\, (\ff,b_2),\, c_2)}{\sigma}
       \toc \cfg{((b_1,b_1),\, (\mathtt{skip},c_1),\, b_2,\, (c_2,c_2))}{\sigma}} \rname{LoopExp2F}
\]}
{\small
\[
\dfrac{}{\cfg{((\ff,b_1),\, (\mathtt{skip},c_1),\, b_2,\, (\mathtt{skip},c_2))}{\sigma}
       \toc \cfg{(b_1,\, (c_1,c_1),\, (b_2,b_2),\, c_2)}{\sigma}} \rname{LoopExp1F}
\]}
\[
\dfrac{}{\cfg{((\ff,b_1),\, c_1,\, b_2,\, c_2)}{\sigma} \toc \bot} \rnameerr{LoopError1}
\]
\[
\dfrac{}{\cfg{((\ttt,b_1),\, (\mathtt{skip},c_1),\, b_2,\, (\mathtt{skip},c_2))}{\sigma}
       \toc \bot} \rnameerr{LoopError2}
\]

\noindent
Il flusso è dunque: \textsc{LoopMainS} entra nel ciclo e arma la valutazione
di $b_1$; \textsc{CtxExp} la porta a un valore, e \textsc{LoopExp1T} (che
richiede $b_1$ \emph{vera}) avvia $c_1$ e arma $b_2$. L'esecuzione di $c_1$
avviene tramite \textsc{CtxStmt}; quando $c_1$ è ridotto a $\mathtt{skip}$,
\textsc{CtxExp} valuta $b_2$: se è vera \textsc{LoopEndS} termina il ciclo in
$\cfg{\mathtt{skip}}{\sigma}$, se è falsa \textsc{LoopExp2F} avvia il
\emph{loop}-corpo $c_2$ e riarma $b_1$. Terminato $c_2$, $b_1$ viene
rivalutata e deve ora essere \emph{falsa}: \textsc{LoopExp1F} riavvia
l'iterazione. Le due regole d'errore coprono i due casi in cui l'asserzione
$b_1$ ha il valore sbagliato: \textcolor{red!60!black}{\textsc{LoopError1}}
quando $b_1$ è falsa \emph{all'ingresso}, \textcolor{red!60!black}{\textsc{LoopError2}}
quando $b_1$ ridiventa \emph{vera} a metà ciclo. Sono le uniche due regole
d'errore del ciclo: ogni altro fallimento nel corpo risale per
\textcolor{red!60!black}{\textsc{CtxError}}.

\paragraph{Il ciclo di Janus.} Kairos adotta \emph{tal quale} il ciclo di
Janus, nella forma e con le regole date dalla base~\cite{LamiLaneseStefani2024},
a sua volta fedele all'originale di Yokoyama e Glück~\cite{YokoyamaGluck2007}.
La forma a un solo corpo, $\mathtt{from}\ b_1\ \mathtt{loop}\ c\
\mathtt{until}\ b_2$, ne è il caso particolare
$c_2 = \mathtt{skip}$, cioè $\mathtt{from}\ b_1\ \mathtt{do}\ c\ \mathtt{loop}\
\mathtt{skip}\ \mathtt{until}\ b_2$; con il secondo corpo si recupera l'idioma
standard di Janus (il lavoro nel \emph{do}-corpo, l'avanzamento del contatore
nel \emph{loop}-corpo), che con un corpo solo non era esprimibile
direttamente.

\subsubsection{Chiamata e inversa}
\label{sec:call}
Sia $\mathtt{procedure}\ f(\bar p)\ \{\, c_f \,\}$. La chiamata sostituisce i
parametri formali $\bar p$ con gli attuali $\bar v$ (lo $\toc$ opera sulle celle
del chiamante). $\mathtt{uncall}$ esegue il corpo \emph{invertito}
$\Inv{c_f}$ (operatore $\Inv{\cdot}$, \S\ref{sec:rev}).
\[
\dfrac{}{\cfg{\mathtt{call}\ f(\bar v)}{\sigma} \toc \cfg{c_f[\bar v/\bar p]}{\sigma}} \rname{Call}
\]
\[
\dfrac{}{\cfg{\mathtt{uncall}\ f(\bar v)}{\sigma} \toc \cfg{\Inv{c_f}[\bar v/\bar p]}{\sigma}} \rname{Uncall}
\]

\subsubsection{Stack}
\label{sec:stack}
$\mathtt{push}(v,s)$ sposta il valore di $v$ in cima a $s$ e \textbf{azzera} $v$;
$\mathtt{pop}(v,s)$ fa il duale (estrae la cima e la mette in $v$), con la
premessa \emph{zero-cleared}: il bersaglio deve essere già $0$.
\[
\dfrac{\sigma(v) = n}
      {\cfg{\mathtt{push}(v,s)}{\sigma} \toc \cfg{\mathtt{skip}}{\upd{\upd{\sigma}{v}{0}}{s}{n : \sigma(s)}}} \rname{Push}
\]
\[
\dfrac{\sigma(v) = 0 \quad \sigma(s) = n : \ell}
      {\cfg{\mathtt{pop}(v,s)}{\sigma} \toc \cfg{\mathtt{skip}}{\upd{\upd{\sigma}{v}{n}}{s}{\ell}}} \rname{Pop}
\]
\[
\dfrac{\sigma(v) = 0 \quad \sigma(s) = \varepsilon}
      {\cfg{\mathtt{pop}(v,s)}{\sigma} \toc \err} \rnameerr{Pop-Err}
\]
\[
\dfrac{\sigma(v) \neq 0}
      {\cfg{\mathtt{pop}(v,s)}{\sigma} \toc \err} \rnameerr{Pop-Err2}
\]
Se $\sigma(v) \neq 0$, \textcolor{red!60!black}{\textsc{Pop-Err2}} dà $\err$.

\subsubsection{Canali: \texttt{ssend} / \texttt{srecv}}
\label{sec:chan}
Un canale $\mathit{ch}$ \textbf{non} contiene valori: è un nome di
sincronizzazione, un punto di rendez-vous tra i due rami che in quel momento
ne sono titolari (vincolo di buona formazione, sotto). $\mathtt{ssend}$ e
$\mathtt{srecv}$ non riducono in isolamento: ciascuno produce un passo
\textbf{etichettato}, che si combina con quello del ramo partner nella
regola \textsc{Par-Comm} (\S\ref{sec:par}) in un unico passo $\toc$ del
\texttt{par} che li contiene, le etichette sono
\[
\alpha ::= \mathit{ch}!\langle \bar v,\bar n\rangle \mid \mathit{ch}?\langle \bar w\rangle
\]
($\mathit{ch}!\langle \bar v,\bar n\rangle$: offerta in invio, variabili
$\bar v$ e valori correnti $\bar n$; $\mathit{ch}?\langle \bar w\rangle$:
offerta in ricezione, variabili bersaglio $\bar w$). I passi etichettati non
modificano da soli lo store, il movimento avviene tutto in
\textsc{Par-Comm}:
\[
\dfrac{\sigma(v_1) = n_1 \quad \cdots \quad \sigma(v_k) = n_k}
      {\cfg{\mathtt{ssend}(\langle \bar v \rangle, \mathit{ch})}{\sigma}
       \xrightarrow{\ \mathit{ch}!\langle \bar v,\bar n\rangle\ } \cfg{\mathtt{skip}}{\sigma}} \rname{Ssend-Lbl}
\]
\[
\dfrac{\sigma(w_1) = 0 \quad \cdots \quad \sigma(w_k) = 0}
      {\cfg{\mathtt{srecv}(\langle \bar w \rangle, \mathit{ch})}{\sigma}
       \xrightarrow{\ \mathit{ch}?\langle \bar w\rangle\ } \cfg{\mathtt{skip}}{\sigma}} \rname{Srecv-Lbl}
\]
\[
\dfrac{\exists i.\ \sigma(w_i) \neq 0}
      {\cfg{\mathtt{srecv}(\langle \bar w \rangle, \mathit{ch})}{\sigma} \toc \err} \rnameerr{Srecv-Err}
\]
\textcolor{red!60!black}{\textsc{Srecv-Err}} non è un passo etichettato: è una premessa su stato
\textbf{locale} al ramo ricevente ($\sigma(\bar w) \neq \bar 0$), quindi
produce $\err$ direttamente, senza
attendere né dipendere da alcun partner, esattamente come
\textcolor{red!60!black}{\textsc{Pop-Err2}}. \textsc{Ssend-Lbl} non ha un errore corrispondente:
leggere $\sigma(\bar v)$ non ha precondizioni, può assumere qualunque valore.

\medskip
\noindent
Il blocco è ora \textbf{emergente}, non il fallimento di una premessa
condivisa: se nessun ramo concorrente offre l'etichetta duale, nessuna
regola dell'intero sistema si applica a quella configurazione: non c'è una
regola \textsc{Ssend}/\textsc{Srecv} "in isolamento" da cui una premessa
possa mancare; semplicemente \textsc{Par-Comm} (\S\ref{sec:par}) non trova
un match. La sincronizzazione stessa \emph{è} la premessa condivisa.

\paragraph{Perché non serve una coda.} Un canale ha al più due titolari
(vincolo di sessione binaria, sotto): in ogni istante c'è al più un
$\mathtt{ssend}$ e un $\mathtt{srecv}$ in attesa sullo stesso canale, e
\textsc{Par-Comm} li consuma insieme nello stesso passo. Non c'è mai più di
un messaggio "in transito" da bufferizzare; se il canale è dentro un ciclo,
questo si ripete a ogni iterazione, ma un invio/ricezione alla volta.

\medskip
\noindent
La ragione \emph{decisiva}, però, non è il conteggio dei messaggi in transito:
è l'\textbf{invertibilità}. Se un canale fosse una coda FIFO,
$\mathtt{ssend}$ e $\mathtt{srecv}$ non potrebbero essere l'uno l'inverso
dell'altro, perché agirebbero su \textbf{estremità opposte}: l'inverso di un
accodamento in fondo è un'estrazione dal fondo, non la $\mathtt{srecv}$ che
preleva dalla testa. Non è l'ordine in sé a rompere la simmetria: si veda
$\mathtt{push}/\mathtt{pop}$ (\S\ref{sec:stack}), che operano su una struttura
ordinata e restano una coppia inversa esatta proprio perché toccano
\emph{entrambi} la cima. È l'asimmetria delle due estremità. Con un unico
valore scambiato in rendez-vous il problema non si pone, e la dualità
$\Inv{\mathtt{ssend}} = \mathtt{srecv}$ enunciata in \S\ref{sec:rev} vale
esattamente, per la stessa convenzione \emph{zero-cleared} che regola
$\mathtt{push}/\mathtt{pop}$. La rinuncia alla coda non è quindi una
semplificazione di comodo, ma ciò che rende i due comandi una coppia inversa.

\paragraph{Due garanzie, e due momenti in cui verificarle.} Su un canale servono
due cose distinte: \emph{chi} può usarlo, e \emph{in che ordine} i due titolari
lo usano. Nella teoria dei tipi di sessione le dà entrambe un unico sistema di
tipi, statico. Qui le due sono separate, e non per comodità di esposizione: la
prima riguarda una titolarità che \emph{cambia durante l'esecuzione}, perché un
canale si può passare, e nessuna analisi del testo può seguirla; la seconda è
invece una proprietà del testo del programma, e si controlla prima di eseguirlo.
Da qui i due colori: la titolarità è \dynamic{sorvegliata a runtime}, il
protocollo \static{verificato in compilazione}.

\paragraph{\dynamic{Sessione binaria, come proprietà dinamica.}} Un canale ha,
in ogni istante, \textbf{al più due} titolari: è una \emph{sessione binaria}.
Non è richiesto un ruolo fisso (mittente o ricevente): i due titolari possono
scambiarsi i ruoli nel tempo, perché con due soli titolari è sempre chiaro chi
sia il partner di una data offerta $\mathit{ch}!$/$\mathit{ch}?$. È ciò che
rende unico l'accoppiamento all'indietro (\S\ref{sec:rev}). Nella letteratura
sui calcoli di processo è la condizione di linearità garantita dai \emph{session
type} binari~\cite{HondaVasconcelosKubo1998}, di cui esiste la generalizzazione
a più di due partecipanti (\emph{multiparty session
types}~\cite{HondaYoshidaCarbone2008}).

\medskip
\noindent
La titolarità è \textbf{stato d'esecuzione}, non una proprietà del testo del
programma, e può cambiare durante l'esecuzione. Un ramo diventa titolare quando
usa il canale o quando lo riceve come payload di un altro canale, e smette di
esserlo quando lo spedisce a un terzo. Da questa sola regola discendono, senza
casi particolari, i due modi in cui un canale si passa:

\begin{itemize}
\item \emph{name passing}, nello stile del $\pi$-calcolo: un ramo crea un canale
privato, lo spedisce e continua a usarlo per ricevere la risposta. Cede la
titolarità spedendolo e la riprende al primo uso successivo, così i titolari
finiscono per essere mittente e destinatario.
\item \emph{delega}: il canale è fra $A$ e $B$, $A$ lo passa a $C$ e smette di
usarlo; i titolari diventano $B$ e $C$. Se $A$ lo riusasse sarebbe un terzo
titolare, e viene fermato. La condizione classica della delega, \emph{chi cede
smette di usarlo prima che il nuovo cominci}, non è qui un obbligo dichiarato a
parte: è una conseguenza del contare al più due titolari alla volta.
\end{itemize}

\noindent
La differenza rispetto alla forma \emph{statica} del vincolo, in cui un canale
appartiene agli stessi due rami per tutta la vita, non è di comodo. Leggendo il
testo del programma si può contare quanti rami \emph{nominano} un canale, ma non
si può seguire un canale che cambia titolare, perché il nuovo titolare lo riceve
sotto un altro nome, magari in un'altra procedura; l'analisi statica è costretta
o a rifiutare la delega o a esentarla dal controllo. Osservando l'esecuzione la
difficoltà non si presenta: l'identità di un canale è il canale stesso, non il
nome della variabile che lo contiene.

\medskip
\noindent
Il vincolo non fa parte della relazione di transizione: è una condizione
sull'esecuzione, la cui violazione produce $\err$. Sono due i modi di violarla, e
corrispondono ai due errori di \S\ref{sec:err}: un terzo ramo che tocca un canale
già impegnato, e lo \emph{stallo}, cioè un ramo che attende un partner che non
può più arrivare.

\paragraph{\static{Tipi di sessione.}} Il vincolo appena dato dice \emph{quanti}
rami toccano un canale, non \emph{come} lo usano: due rami che eseguono entrambi
un $\mathtt{ssend}$ lo soddisfano, e restano bloccati per sempre. Un
\textbf{tipo di sessione} descrive la sequenza degli scambi che un ramo compie
su un canale:
\[
S ::= \mathtt{end} \mid\ !k.S \mid\ ?k.S \mid \mu X.S \mid X
\]
dove $k$ è il numero di interi scambiati in un passo. La \textbf{dualità}
scambia le due direzioni e lascia il resto:
\[
\overline{\mathtt{end}} = \mathtt{end} \qquad
\overline{!k.S} =\ ?k.\overline{S} \qquad
\overline{?k.S} =\ !k.\overline{S} \qquad
\overline{\mu X.S} = \mu X.\overline{S} \qquad
\overline{X} = X
\]
Il tipo di un ramo si ricostruisce dal programma: $!k$ per un $\mathtt{ssend}$ di
$k$ valori, $?k$ per un $\mathtt{srecv}$, concatenazione per la sequenza,
$\mu X.S.X$ per un ciclo il cui corpo ha tipo $S$, e il tipo del corpo per una
chiamata di procedura, con il canale sostituito al parametro formale. Due rami
usano bene un canale quando i loro tipi sono duali, e la dualità è
\textbf{invariante per inversione}: invertire un \texttt{par} scambia
$\mathtt{ssend}$ e $\mathtt{srecv}$ in entrambi i rami (\S\ref{sec:rev}), cioè
dualizza entrambi i tipi, e due tipi duali restano duali.

\medskip
\noindent
Questo controllo non sostituisce quello dinamico e non lo duplica: riconosce in
anticipo i conflitti già visibili nel testo, come due rami che offrono la stessa
direzione o si scambiano quantità diverse di valori, che altrimenti si
manifesterebbero solo come stallo a esecuzione avviata. Il frammento è minimo e va preso per quello
che è: non c'è \emph{scelta etichettata}, quindi i due rami di un $\mathtt{if}$
che comunichi devono compiere le stesse azioni sul canale, e non ci sono
sessioni \emph{multiparty}. Dove il frammento non basta a descrivere l'uso di un
canale, il tipo resta indeterminato e non si segnala nulla: si riconoscono i soli
conflitti dimostrabili, mai rifiutando un programma corretto per insufficienza
della descrizione.

\subsubsection{Composizione parallela: \texttt{par\,\dots\,and\,\dots\,rap}}
\label{sec:par}
I rami procedono per \emph{interleaving} non deterministico (caso a due rami; con
più \texttt{and} si generalizza). Un passo avanza il ramo sinistro o il destro:
\[
\dfrac{\cfg{c_1}{\sigma} \toc \cfg{c_1'}{\sigma'}}
      {\cfg{\mathtt{par}\ c_1\ \mathtt{and}\ c_2\ \mathtt{rap}}{\sigma} \toc
       \cfg{\mathtt{par}\ c_1'\ \mathtt{and}\ c_2\ \mathtt{rap}}{\sigma'}} \rname{Par-L}
\]
\[
\dfrac{\cfg{c_2}{\sigma} \toc \cfg{c_2'}{\sigma'}}
      {\cfg{\mathtt{par}\ c_1\ \mathtt{and}\ c_2\ \mathtt{rap}}{\sigma} \toc
       \cfg{\mathtt{par}\ c_1\ \mathtt{and}\ c_2'\ \mathtt{rap}}{\sigma'}} \rname{Par-R}
\]
Quando un ramo termina (cioè è diventato $\mathtt{skip}$), lo si elimina
dal blocco:
\[
\dfrac{}{\cfg{\mathtt{par}\ \mathtt{skip}\ \mathtt{and}\ c_2\ \mathtt{rap}}{\sigma} \toc
       \cfg{c_2}{\sigma}} \rname{Par-L0}
\]
\[
\dfrac{}{\cfg{\mathtt{par}\ c_1\ \mathtt{and}\ \mathtt{skip}\ \mathtt{rap}}{\sigma} \toc
       \cfg{c_1}{\sigma}} \rname{Par-R0}
\]
(Quando resta un solo ramo, $\mathtt{par}\ c\ \mathtt{rap}$ si comporta come
$c$.) \textsc{Par-L} e \textsc{Par-R} sono \textsc{CtxStmt} (\S\ref{sec:base})
istanziata sulle due clausole $C_s[\bullet] = \mathtt{par}\ \bullet\
\mathtt{and}\ c\ \mathtt{rap}$ e $\mathtt{par}\ c\ \mathtt{and}\ \bullet\
\mathtt{rap}$: le riportiamo per esteso perché è qui che si vede
l'interleaving, ma non sono regole indipendenti. \textsc{Par-L0}/\textsc{Par-R0}
lo sono, perché riguardano la terminazione di un ramo. Se un ramo fallisce,
l'intero blocco fallisce (un errore non è localizzabile a un solo ramo una
volta che i rami condividono il blocco \texttt{par}), ma anche questo non
richiede regole proprie: il fallimento di un ramo risale all'intero blocco
per \textcolor{red!60!black}{\textsc{CtxError}} sulle stesse due clausole
(\S\ref{sec:err}).

\medskip
\noindent
\textsc{Par-L}/\textsc{Par-R} (e le varianti \textsc{-L0}/\textsc{-R0})
sollevano un passo $\toc$ \textbf{ordinario} (non etichettato) di un solo
ramo. Un $\mathtt{ssend}$/$\mathtt{srecv}$ non produce mai un passo $\toc$ da
solo (\S\ref{sec:chan}), quindi queste regole non lo sollevano mai
direttamente; ma se il ramo che lo contiene è a sua volta un \texttt{par}
annidato, l'offerta deve poter risalire fino al livello dove si trova il suo
partner. Due regole gemelle fanno da tramite, sollevando un passo
\textbf{etichettato} anziché ordinario:
\[
\dfrac{\cfg{c_1}{\sigma} \xrightarrow{\ \alpha\ } \cfg{c_1'}{\sigma}}
      {\cfg{\mathtt{par}\ c_1\ \mathtt{and}\ c_2\ \mathtt{rap}}{\sigma}
       \xrightarrow{\ \alpha\ } \cfg{\mathtt{par}\ c_1'\ \mathtt{and}\ c_2\ \mathtt{rap}}{\sigma}} \rname{Par-L-Lbl}
\]
\[
\dfrac{\cfg{c_2}{\sigma} \xrightarrow{\ \alpha\ } \cfg{c_2'}{\sigma}}
      {\cfg{\mathtt{par}\ c_1\ \mathtt{and}\ c_2\ \mathtt{rap}}{\sigma}
       \xrightarrow{\ \alpha\ } \cfg{\mathtt{par}\ c_1\ \mathtt{and}\ c_2'\ \mathtt{rap}}{\sigma}} \rname{Par-R-Lbl}
\]
Con queste, un'offerta di canale generata da \textsc{Ssend-Lbl}/
\textsc{Srecv-Lbl} in fondo a un \texttt{par} annidato risale, un livello di
nesting alla volta, fino al primo livello in cui trova un ramo gemello con
l'offerta duale sullo stesso canale, esattamente lo stile standard alla
CCS. La sincronizzazione vera e propria resta un'unica regola, applicata al
livello dove le due offerte sono gemelle:
\[
\dfrac{\cfg{c_1}{\sigma} \xrightarrow{\ \mathit{ch}!\langle \bar v,\bar n\rangle\ } \cfg{c_1'}{\sigma}
       \qquad
       \cfg{c_2}{\sigma} \xrightarrow{\ \mathit{ch}?\langle \bar w\rangle\ } \cfg{c_2'}{\sigma}}
      {\cfg{\mathtt{par}\ c_1\ \mathtt{and}\ c_2\ \mathtt{rap}}{\sigma} \toc
       \cfg{\mathtt{par}\ c_1'\ \mathtt{and}\ c_2'\ \mathtt{rap}}{\upd{\upd{\sigma}{\bar v}{\bar 0}}{\bar w}{\bar n}}}
      \rname{Par-Comm}
\]
(simmetrica in $c_1,c_2$: si applica anche scambiando quale dei due offre
$\mathit{ch}!$ e quale $\mathit{ch}?$; con più di due \texttt{and} si applica
alla coppia di rami che offre le etichette duali). Il canale $\mathit{ch}$
compare solo per far combaciare le due etichette: \textbf{non} compare
nell'aggiornamento dello store, che tocca solo $\bar v$ (azzerate) e $\bar w$
(riempite con $\bar n$), è il passo congiunto descritto in
\S\ref{sec:chan}. \textsc{Par-Comm} produce un passo $\toc$ ordinario (non
etichettato): una volta sincronizzato, il risultato non deve risalire
oltre, è \emph{questo} passo, non l'offerta grezza, che \textsc{Par-L}/
\textsc{Par-R} sollevano ai livelli ancora più esterni.

\paragraph{\textsc{Par-L-Lbl}/\textsc{Par-R-Lbl} non aprono un instradamento
sbagliato.} Sollevare un'etichetta attraverso un \texttt{par} annidato
potrebbe in linea di principio farla incontrare un partner diverso da quello
inteso, se sullo stesso canale ci fossero più di due offerte possibili in
giro per l'albero dei \texttt{par}. Ma il vincolo di sessione binaria
(\S\ref{sec:chan}), \dynamic{sorvegliato durante l'esecuzione}, garantisce
che in ogni istante siano \textbf{al più due} i rami (a qualunque profondità di
nesting) titolari di un dato canale, e quindi al più due quelli che possono
offrire un'etichetta su di esso:
\textsc{Par-Comm} confronta le etichette per \emph{nome di canale}, quindi se
solo due rami possono mai offrirne una su $\mathit{ch}$, non c'è un terzo
candidato con cui una $\mathtt{Par}$-\textsc{Lbl} possa accoppiarla per
errore, qualunque cammino di risalita la porti a incontrare l'altra offerta
su $\mathit{ch}$, è per forza quella giusta. Il controllo statico
(implementato in \texttt{parser.py}, non riportato qui) conta esattamente
questo: le foglie concorrenti (i rami che restano dopo aver scomposto
ricorsivamente ogni \texttt{par} annidato nei suoi stessi rami, non i soli
rami immediati) che toccano ciascun canale, indipendentemente da quanto in
profondità si trovino; è la stessa nozione di "ramo" che qui risale con
\textsc{Par-L-Lbl}/\textsc{Par-R-Lbl}.

\medskip
\noindent
Con \texttt{par} la relazione $\toc$ \textbf{non} è più deterministica: l'ordine
d'interleaving non è fissato. Tuttavia i rami \textbf{non condividono memoria}:
ciascuno opera su una porzione \emph{disgiunta} dello store (variabili
\texttt{int}/\texttt{stack} proprie). L'\textbf{unica} forma di
comunicazione/sincronizzazione tra rami passa per i canali, tramite
\textsc{Par-Comm}. Grazie a questa disgiunzione, i diversi interleaving che non
toccano canali producono lo stesso store finale (confluenza).

\paragraph{Perché vale la confluenza.}
\label{sec:par-conf}
Siano $\cfg{\mathtt{par}\ c_1\
\mathtt{and}\ c_2\ \mathtt{rap}}{\sigma} \toc \cfg{\mathtt{par}\ c_1'\
\mathtt{and}\ c_2\ \mathtt{rap}}{\sigma_1}$ (via \textsc{Par-L}) e
$\cfg{\mathtt{par}\ c_1\ \mathtt{and}\ c_2\ \mathtt{rap}}{\sigma} \toc
\cfg{\mathtt{par}\ c_1\ \mathtt{and}\ c_2'\ \mathtt{rap}}{\sigma_2}$ (via
\textsc{Par-R}) i due passi abilitati in $\sigma$, nessuno dei due un
$\mathtt{ssend}/\mathtt{srecv}$ (nessun tocco a canali). Poiché $c_1$ legge e
scrive solo le proprie variabili \texttt{int}/\texttt{stack} e $c_2$ solo le
proprie, e questi due insiemi sono disgiunti, l'aggiornamento di $\sigma$
compiuto dal passo di $c_1$ e quello compiuto dal passo di $c_2$ agiscono su
porzioni disgiunte del dominio: sono \emph{indipendenti} e commutano.
Applicando prima l'uno e poi l'altro (in un ordine o nell'altro) si ottiene
quindi lo stesso $\sigma_{12}$, con $\cfg{\mathtt{par}\ c_1'\ \mathtt{and}\
c_2\ \mathtt{rap}}{\sigma_1} \toc \cfg{\mathtt{par}\ c_1'\ \mathtt{and}\
c_2'\ \mathtt{rap}}{\sigma_{12}}$ e simmetricamente da $\sigma_2$: un
diamante di confluenza locale. Estendendo per induzione sulla lunghezza della
derivazione (ogni permutazione di passi indipendenti è una sequenza di
scambi adiacenti di questo tipo) si ottiene che due interleaving qualunque
che non toccano canali portano allo stesso store finale. La clausola è
necessaria: un passo \textsc{Par-Comm} coinvolge \textbf{entrambi} i rami
contemporaneamente (\S\ref{sec:par}), non un passo di solo $c_1$ sollevato da
\textsc{Par-L} indipendentemente da un passo di solo $c_2$, non c'è quindi
un passo-di-$c_1$ e un passo-di-$c_2$ separati da far commutare, ce n'è uno
solo che li consuma entrambi in un colpo; non rientra nel caso trattato qui.
(\S\ref{sec:rev} tratta l'argomento, analogo ma non lo stesso, per
l'accoppiamento forzato all'indietro su un canale.)

\subsubsection{Commit / rollback: \texttt{try\,\dots\,rollback\,\dots\,yrt}}
\label{sec:try}
Esegue il corpo $c$ in avanti, poi valuta la guardia di commit $b'$ sullo store
risultante: se $b'$ è \emph{vera} il corpo resta applicato (commit); se è
\emph{falsa} il corpo viene \textbf{annullato} eseguendo l'inverso $\Inv{c}$
(\S\ref{sec:rev}), che riporta esattamente allo store iniziale, e poi si esegue
il corpo di rollback $r$. A differenza di \texttt{if\,\dots\,fi}, le due
guardie $b$ (ingresso) e $b'$ (uscita) non sono legate da alcun vincolo di
coerenza: la decisione commit/rollback dipende solo da $b'$ valutata a fine
corpo. La guardia d'ingresso $b$ ha quindi valore puramente documentale: non
compare in alcuna premessa, né nello zucchero sintattico. Darle un ruolo
semantico (una regola d'ingresso che la valuti, sul modello dell'asserzione
del ciclo) è una possibile estensione, che qui non perseguiamo.

\medskip
\noindent
Nella terminologia standard il costrutto è una \textbf{transazione}: con corpo
di rollback vuoto ($r = \mathtt{skip}$) è una transazione \emph{acida} (o si
applica interamente o non lascia traccia), mentre con un corpo di rollback non
vuoto è una transazione \textbf{con compensazione}, dove $r$ è appunto
l'azione compensativa eseguita dopo aver disfatto il tentativo. Il rollback in
contesto concorrente è studiato in~\cite{LaneseMezzinaSchmittStefani2011,
LaneseLienhardtMezzinaSchmittStefani2013}. Diamo qui le regole in stile
big-step.

\[
\dfrac{\cfg{c}{\sigma} \toc^{*} \cfg{\mathtt{skip}}{\sigma'} \quad \cfg{b'}{\sigma'} \tob \cfg{\ttt}{\sigma'}}
      {\cfg{\mathtt{try}\ b\ c\ \mathtt{rollback}\ r\ \mathtt{yrt}\ b'}{\sigma} \toc \cfg{\mathtt{skip}}{\sigma'}}
      \rname{Try-Commit}
\]
\[
\dfrac{\begin{array}{c}
       \cfg{c}{\sigma} \toc^{*} \cfg{\mathtt{skip}}{\sigma'} \quad \cfg{b'}{\sigma'} \tob \cfg{\ff}{\sigma'}\\[0.3em]
       \cfg{\Inv{c}}{\sigma'} \toc^{*} \cfg{\mathtt{skip}}{\sigma} \quad \cfg{r}{\sigma} \toc^{*} \cfg{\mathtt{skip}}{\sigma''}
       \end{array}}
      {\cfg{\mathtt{try}\ b\ c\ \mathtt{rollback}\ r\ \mathtt{yrt}\ b'}{\sigma} \toc \cfg{\mathtt{skip}}{\sigma''}}
      \rname{Try-Rollback}
\]
\[
\dfrac{\cfg{c}{\sigma} \toc^{*} \err}
      {\cfg{\mathtt{try}\ b\ c\ \mathtt{rollback}\ r\ \mathtt{yrt}\ b'}{\sigma} \toc \err}
      \rnameerr{Try-Err}
\]
La premessa $\cfg{\Inv{c}}{\sigma'} \toc^{*} \sigma$ è garantita dalla proprietà
di reversibilità (\S\ref{sec:rev}): poiché $\cfg{c}{\sigma} \toc^{*} \sigma'$,
l'inverso del corpo riporta allo store $\sigma$ \emph{pre-corpo}. Senza clausola
\texttt{rollback} si pone $r = \mathtt{skip}$ (solo annullamento). È il
\textbf{backtracking nativo}: il ripristino dello stato è l'inverso esatto del
tentativo, senza copie/snapshot. \textcolor{red!60!black}{\textsc{Try-Err}} lascia che un errore del
corpo esca dal \texttt{try} inalterato: un $\err$ è un bug di programma, non
un fallimento della guardia $b'$, quindi non attiva il rollback.

\paragraph{\static{Vincolo di buona formazione: niente canali nel corpo.}} Le premesse
$\cfg{c}{\sigma} \toc^{*} \sigma'$ e $\cfg{\Inv{c}}{\sigma'} \toc^{*} \sigma$
usano la relazione \textbf{non} etichettata: un $\mathtt{ssend}$/$\mathtt{srecv}$
dentro $c$ produce però solo un passo etichettato (\S\ref{sec:chan}), e
l'unica regola che lo consuma, \textsc{Par-Comm} (\S\ref{sec:par}), vive al
livello di un \texttt{par}, sopra $c$, non dentro. Se $c$ (o $r$) contiene
un $\mathtt{ssend}$/$\mathtt{srecv}$, diretto o raggiunto tramite
$\mathtt{call}$/$\mathtt{uncall}$, $\cfg{c}{\sigma} \toc^{*} \sigma'$ non è
derivabile: il \texttt{try} resta bloccato, non fallisce con $\err$. Un
programma ben formato evita quindi \texttt{ssend}/\texttt{srecv} nel corpo e
nel rollback di ogni \texttt{try}, direttamente o tramite
\texttt{call}/\texttt{uncall} (\static{verificato staticamente}, non nella relazione
di transizione). Non è un limite che valga la pena aggirare: anche
potendolo esprimere, disfare uno $\mathtt{srecv}$ già sincronizzato
richiederebbe di far ripartire dal punto giusto anche il ramo mittente, che
nel frattempo può aver già proseguito, un \texttt{try}/\texttt{rollback} è
un costrutto a un solo ramo, e non ha modo di coordinare la cooperazione di
un partner che non controlla. Non è una limitazione ad hoc: è la difficoltà
nota del rollback in contesto concorrente, dove annullare un'azione già
comunicata obbliga a propagare l'annullamento ai partner che l'hanno
osservata~\cite{LaneseMezzinaSchmittStefani2011,
LaneseLienhardtMezzinaSchmittStefani2013}.

\noindent
Realizzazione (zucchero sintattico): il costrutto si riscrive con \texttt{if}
e \texttt{call}/\texttt{uncall} su una procedura sintetica $f_c$ con corpo $c$
(parametri = variabili libere di $c$):
\[
\mathtt{call}\ f_c(\bar v)\,;\,
\mathtt{if}\ b'\ \mathtt{then}\ \mathtt{skip}\ \mathtt{else}\
   \mathtt{uncall}\ f_c(\bar v)\,;\, r\ \mathtt{fi}\ b' .
\]

% ======================================================================
\section{Errori a runtime}
\label{sec:err}
% ======================================================================

L'esito d'errore è $\err$, ereditato dalla base insieme alla regola che lo
solleva: $\Gamma_c' = \Gamma_c \cup \{\err\}$, $\top_c' = \top_c \cup
\{\err\}$. La \textbf{propagazione} è interamente a carico della base:
\textcolor{red!60!black}{\textsc{CtxError}} (\S\ref{sec:base}) fa risalire un
$\err$ da qualunque posizione di contesto $C_s[\bullet]$: attraverso una
sequenza, i rami di un condizionale, i corpi di un ciclo e i rami di un
\texttt{par}. Nessun costrutto ha quindi bisogno di una propria
regola di propagazione. Resta un'eccezione, \textcolor{red!60!black}{\textsc{Try-Err}}
(\S\ref{sec:try}): il corpo di un \texttt{try} non è un buco di contesto,
perché le sue regole lo riducono con $\toc^{*}$ in una premessa e devono
conservarlo intatto per poterne eseguire l'inverso $\Inv{c}$ in caso di
rollback. Lì una regola d'errore esplicita serve ancora.

\medskip
\noindent
Quel che resta di specifico a Kairos è \textbf{quali} costrutti possono
fallire \emph{localmente}. Sono i costrutti che Kairos aggiunge alla base
(\S\ref{sec:base}), e ciascuno ha una regola propria che produce $\err$ invece
di restare bloccato, perché nessun altro ramo potrebbe mai soddisfare quel
vincolo, trattandosi di stato \textbf{locale} al ramo (\S\ref{sec:par}):
\textcolor{red!60!black}{\textsc{Local-Err}} ($v \in \mathrm{dom}(\sigma)$) e
\textcolor{red!60!black}{\textsc{Deloc-Err}} ($v \notin \mathrm{dom}(\sigma)$
oppure $\sigma(v) \neq m$) per lo scope (\S\ref{sec:scope});
\textcolor{red!60!black}{\textsc{Pop-Err}} (stack vuoto) e
\textcolor{red!60!black}{\textsc{Pop-Err2}} ($\sigma(v) \neq 0$) per lo stack
(\S\ref{sec:stack}); \textcolor{red!60!black}{\textsc{Srecv-Err}}
($\sigma(\bar w) \neq \bar 0$) per la ricezione su canale (\S\ref{sec:chan}).
Gli errori dei costrutti \emph{ereditati}, cioè l'asserzione incoerente nel
condizionale (\textcolor{red!60!black}{\textsc{IfError1}}/%
\textcolor{red!60!black}{\textsc{IfError2}}) e nel ciclo
(\textcolor{red!60!black}{\textsc{LoopError1}}/%
\textcolor{red!60!black}{\textsc{LoopError2}}), vengono invece dalla base e
non sono discussi qui.

\medskip
\noindent
\textbf{$\err$ contro blocco.} Questa distinzione è propria di Kairos, perché
è la concorrenza a renderla possibile. Il blocco di sincronizzazione sui
canali (\S\ref{sec:chan}) non è il fallimento di una premessa: è l'assenza di
un ramo partner pronto con cui combinarsi in \textsc{Par-Comm}
(\S\ref{sec:par}): nessuna regola si applica finché il partner non è pronto,
ma un ramo parallelo può ancora offrirsi in un momento successivo, quindi
\textbf{blocco}, non $\err$. Il criterio è dunque: se la condizione violata
riguarda stato locale al ramo, nessuna evoluzione dell'ambiente potrà mai
renderla vera e l'esito è $\err$; se riguarda la disponibilità di un partner,
l'esito è blocco. È per questo che
\textcolor{red!60!black}{\textsc{Srecv-Err}} produce $\err$ e non blocco: la
sua premessa $\sigma(\bar w) \neq \bar 0$ parla delle variabili bersaglio del
ricevente, non del canale.

\medskip
\noindent
\textbf{Violazioni della sessione binaria.} Il vincolo di \S\ref{sec:chan} è
una condizione sull'esecuzione, e si viola in due modi, entrambi con esito
$\err$. Il primo è un \textbf{terzo titolare}: un ramo tocca un canale che ne
ha già due. Non è blocco, perché nessuna evoluzione dell'ambiente può liberare
un posto: chi cede un canale lo fa spedendolo, e l'ha già fatto o non lo farà.
Il secondo è lo \textbf{stallo}: tutti i rami di un \texttt{par} sono fermi in
attesa su canale. Il singolo blocco resta blocco, come detto sopra, perché un
partner potrebbe ancora offrirsi; ma quando non resta alcun ramo che possa
offrirsi, l'attesa non è più provvisoria e diventa $\err$. È la differenza fra
non avere \emph{ancora} un partner e non poterne avere.

\medskip
\noindent
La funzione di esecuzione
\[
\mathit{exec} : \mathit{Com} \times \Store \rightharpoonup \Store \cup \{\err\}
\]
resta perciò parziale (un ciclo può divergere, un canale può bloccarsi) ma
segnala esplicitamente le violazioni di reversibilità.

% ======================================================================
\chapter{Reversibilità}
\label{chap:reversibilita}
\label{sec:rev}
% ======================================================================

L'esecuzione all'indietro è data dall'operatore di \emph{inversione}
$\Inv{\cdot}$, definito per induzione sulla struttura del comando:
{\small
\begin{align*}
\Inv{\mathtt{skip}} &= \mathtt{skip} &
\Inv{c_1\,;\,c_2} &= \Inv{c_2}\,;\,\Inv{c_1} \quad \text{(ordine rovesciato)}\\
\Inv{v \mathbin{+\!\!=} e} &= v \mathbin{-\!\!=} e &
\Inv{v \mathbin{-\!\!=} e} &= v \mathbin{+\!\!=} e \\
\Inv{v \mathbin{\hat{}\!=} e} &= v \mathbin{\hat{}\!=} e &
\Inv{v_1 \mathbin{<\!=\!>} v_2} &= v_1 \mathbin{<\!=\!>} v_2 \\
\Inv{\mathtt{local}\ v = m} &= \mathtt{delocal}\ v = m &
\Inv{\mathtt{delocal}\ v = m} &= \mathtt{local}\ v = m \\
\Inv{\mathtt{push}(v,s)} &= \mathtt{pop}(v,s) &
\Inv{\mathtt{pop}(v,s)} &= \mathtt{push}(v,s) \\
\Inv{\mathtt{call}\ f(\bar v)} &= \mathtt{uncall}\ f(\bar v) &
\Inv{\mathtt{uncall}\ f(\bar v)} &= \mathtt{call}\ f(\bar v) \\
\Inv{\mathtt{ssend}(\langle \bar v \rangle, \mathit{ch})} &= \mathtt{srecv}(\langle \bar v \rangle, \mathit{ch}) \\
\Inv{\mathtt{srecv}(\langle \bar v \rangle, \mathit{ch})} &= \mathtt{ssend}(\langle \bar v \rangle, \mathit{ch})
\end{align*}
}
\[
\Inv{\mathtt{if}\ b\ \mathtt{then}\ c_t\ \mathtt{else}\ c_e\ \mathtt{fi}\ b'}
   = \mathtt{if}\ b'\ \mathtt{then}\ \Inv{c_t}\ \mathtt{else}\ \Inv{c_e}\ \mathtt{fi}\ b
\]
\[
\Inv{\mathtt{from}\ b_1\ \mathtt{do}\ c_1\ \mathtt{loop}\ c_2\ \mathtt{until}\ b_2}
   = \begin{array}[t]{@{}l@{}}\mathtt{from}\ b_2\ \mathtt{do}\ \Inv{c_1}\\ \quad\mathtt{loop}\ \Inv{c_2}\ \mathtt{until}\ b_1\end{array}
\]
\[
\Inv{\mathtt{par}\ c_1\ \mathtt{and}\ \dots\ \mathtt{and}\ c_n\ \mathtt{rap}}
   = \mathtt{par}\ \Inv{c_1}\ \mathtt{and}\ \dots\ \mathtt{and}\ \Inv{c_n}\ \mathtt{rap}
\]

\noindent
Le riscritture chiave sono lo \textbf{scambio delle guardie}: nel condizionale
$b \leftrightarrow b'$, nel ciclo $b_1 \leftrightarrow b_2$. È ciò che permette
all'esecuzione inversa di scegliere il ramo giusto / sapere quando fermarsi
senza memorizzare la storia.

\smallskip
\noindent
Il $\mathtt{par}$ si inverte ramo per ramo, senza riordinarli fra loro, e le due
proprietà che il linguaggio già impone bastano a giustificarlo. I rami operano su
porzioni disgiunte dello store (\S\ref{sec:par}), quindi ciò che un ramo fa alle
proprie variabili si inverte indipendentemente da cosa stia facendo un altro; e
ogni canale ha due soli titolari (\S\ref{sec:chan}), quindi ogni appuntamento ha
un unico partner possibile. Poiché $\Inv{\mathtt{ssend}} = \mathtt{srecv}$ e
$\Inv{\mathtt{srecv}} = \mathtt{ssend}$, rovesciando l'ordine dei comandi dentro
ciascun ramo si rovescia la successione degli appuntamenti, e i due estremi di
ciascuno restano l'uno di fronte all'altro. La dimostrazione è più sotto, nel
caso induttivo \texttt{par}.

\smallskip
\noindent
L'argomento ha un'ipotesi, ed è che i titolari di ogni canale siano fissati dal
testo del programma. Se un canale può viaggiare \emph{dentro} un messaggio ---
$\mathtt{ssend}(\langle x, \mathit{ch}_2\rangle, \mathit{ch}_1)$, la mobilità
del $\pi$-calcolo --- la titolarità di $\mathit{ch}_2$ non è più una proprietà
del testo ma il risultato di un appuntamento precedente su $\mathit{ch}_1$.
Rovesciando l'ordine dei comandi, gli usi di $\mathit{ch}_2$ verrebbero prima
dell'appuntamento che lo consegna, e i due estremi non si troverebbero più l'uno
di fronte all'altro: l'esecuzione all'indietro si blocca. Su quei programmi
$\Inv{\cdot}$ non è definita per il $\mathtt{par}$, e l'inversione va ottenuta
per altra via (\S\ref{sec:interprete}). È una restrizione che si riconosce sul
testo, e nel corpus di prova riguarda due programmi su ottantanove.

\smallskip
\noindent
$\mathtt{ssend}$ e $\mathtt{srecv}$ sono, come ogni altra coppia della
tabella, l'una l'inversa dell'altra, senza bisogno di comandi ausiliari.
Questo vale grazie al vincolo di sessione binaria (\S\ref{sec:chan}): con al
più due titolari del canale, la sincronizzazione \textsc{Par-Comm}
(\S\ref{sec:par}) ha sempre un unico partner possibile, in avanti come
all'indietro, quindi scambiare $\mathtt{ssend} \leftrightarrow \mathtt{srecv}$
non può mai accoppiare i passi con il partner sbagliato. Il caso
\texttt{par} più sotto sviluppa questo argomento.

\paragraph{Proprietà di reversibilità.}
Per ogni comando $c$ ben formato e stati $\sigma, \sigma'$ (senza $\err$):
\[
\cfg{c}{\sigma} \toc^{*} \cfg{\mathtt{skip}}{\sigma'}
\qquad\Longleftrightarrow\qquad
\cfg{\Inv{c}}{\sigma'} \toc^{*} \cfg{\mathtt{skip}}{\sigma} .
\]
Cioè $\Inv{c}$ calcola la relazione inversa di $c$ sullo store. In particolare,
escludendo \texttt{par}, $\toc$ è deterministica e $c$ denota una funzione
parziale \emph{iniettiva}; di conseguenza
$\cfg{c\,;\,\Inv{c}}{\sigma} \toc^{*} \sigma$ (identità sullo store).

Il determinismo non è un'osservazione a parte: comando per comando, le
premesse delle sue regole (successo ed $\err$) sono a due a due incompatibili
e la loro unione copre ogni $\sigma$ possibile, sono cioè una
\textbf{partizione} dello spazio degli stati, quindi esattamente una regola
si applica in ogni configurazione. Ad esempio per $\mathtt{pop}$:
$\sigma(v)\neq 0$ (\textcolor{red!60!black}{\textsc{Pop-Err2}}), $\sigma(v)=0 \land \sigma(s)=\varepsilon$
(\textcolor{red!60!black}{\textsc{Pop-Err}}), $\sigma(v)=0 \land \sigma(s)=n : \ell$ (\textsc{Pop})
coprono ogni $\sigma$ esattamente una volta, senza sovrapposizioni. Fa
eccezione solo la sincronizzazione di canale (\S\ref{sec:chan}): lì la
partizione locale resta valida (\textcolor{red!60!black}{\textsc{Srecv-Err}} copre comunque
esattamente il proprio caso), ma \textsc{Par-Comm} può temporaneamente non
trovare un ramo partner pronto, lasciando nessuna regola applicabile: è il
blocco di sincronizzazione, non una
rottura del determinismo: nessuna configurazione ha mai \emph{due} regole
applicabili contemporaneamente.

\medskip
\noindent
La dimostrazione è per induzione strutturale su $c$ (equivalentemente, sulla
derivazione di $\cfg{c}{\sigma} \toc^{*} \sigma'$): nei casi base
gli assegnamenti si annullano per i vincoli $v \notin e$
(\textsc{Add2}$\leftrightarrow$\textsc{Sub2}) o sono involuzioni
($\mathbin{\hat{}\!=}$, $\mathbin{<\!=\!>}$); $\mathtt{push}/\mathtt{pop}$ si
annullano per la convenzione \emph{zero-cleared}: \textsc{Pop} richiede in
ingresso che il bersaglio sia già $0$, e \textsc{Push} lo ristabilisce
esattamente in uscita. $\mathtt{ssend}/\mathtt{srecv}$ (via
\textsc{Par-Comm}, \S\ref{sec:par}) si annullano per la \textbf{stessa}
convenzione, applicata alla coppia di rami invece che a un solo comando: la
premessa $\sigma(\bar w) = \bar 0$ sul lato ricevente ed effetto $\bar
w \mapsto \bar n$ sono zero-cleared esattamente come \textsc{Pop}, e l'effetto
$\bar v \mapsto \bar 0$ sul lato mittente è zero-cleared esattamente come
\textsc{Push}: è ciò che rende $\mathtt{ssend}\,;\,\mathtt{srecv}$
l'identità anche sulle variabili coinvolte, non solo un'astrazione sul
canale;
nei casi induttivi la doppia guardia di \texttt{if}/\texttt{from} garantisce
che l'esecuzione inversa ripercorra lo stesso ramo / le stesse iterazioni.

\smallskip
\noindent
Il caso induttivo \texttt{par}: sia $\cfg{\mathtt{par}\ c_1\ \mathtt{and}\
c_2\ \mathtt{rap}}{\sigma} \toc^{n} \sigma'$ una derivazione di lunghezza $n$.
Ogni passo è o un passo \textsc{Par-L}/\textsc{Par-L0} (risp.
\textsc{Par-R}/\textsc{Par-R0}) che avanza \emph{un solo} ramo di un passo
ordinario locale (esattamente uno dei casi già trattati sopra) o un passo
\textsc{Par-Comm} (\S\ref{sec:par}) che avanza \emph{entrambi} i rami insieme
in una sincronizzazione di canale. Per IH ogni passo locale del primo tipo è
invertito esattamente da $\Inv{\cdot}$ applicato al comando ridotto in quel
passo, indipendentemente da cosa contenga in quel momento l'altro ramo: le
variabili \texttt{int}/\texttt{stack} sono disgiunte tra rami
(\S\ref{sec:par}). Per il secondo tipo, \textsc{Par-Comm} stessa mostra come
si inverte: le sue premesse sono $\sigma(\bar v) = \bar n$ (lato mittente) e
$\sigma(\bar w) = \bar 0$ (lato ricevente), la conclusione è $\sigma[\bar
v \mapsto \bar 0][\bar w \mapsto \bar n]$ (dallo stato risultante,
$\sigma(\bar w) = \bar n \neq \bar 0$ e $\sigma(\bar v) = \bar 0$ sono
esattamente le premesse per rieseguire \textsc{Par-Comm} con mittente e
ricevente scambiati, cioè con $\mathtt{ssend} \leftrightarrow \mathtt{srecv}$
scambiati in ciascun ramo), che riporta esattamente a $\sigma$.

Percorrendo la derivazione all'indietro passo per passo, nell'ordine esatto
inverso a quello forward, e invertendo ad ogni passo ciò che lo ha compiuto
(un ramo solo per \textsc{Par-L}/\textsc{Par-R}, entrambi per
\textsc{Par-Comm}), si ricostruisce $\sigma$: questo è esattamente ciò che
$\Inv{\mathtt{par}\ c_1\ \mathtt{and}\ c_2\ \mathtt{rap}} = \mathtt{par}\
\Inv{c_1}\ \mathtt{and}\ \Inv{c_2}\ \mathtt{rap}$ realizza staticamente.
Resta da giustificare perché l'esecuzione all'indietro trovi \emph{sempre}
questo percorso, e non un altro: non basta che un percorso corretto esista. Con più di due rami candidati su un canale un'esecuzione all'indietro
potrebbe accoppiare i passi con il partner sbagliato, esattamente il tipo di
errore che il vincolo di sessione binaria evita per costruzione. Con
\textbf{esattamente due} rami sul canale (\S\ref{sec:chan}), e purché siano i
due che il testo del programma indica --- è qui che serve l'ipotesi sopra: un
canale che viaggia dentro un messaggio i suoi titolari li acquisisce a tempo di
esecuzione ---, l'accoppiamento è
\textbf{unico}, non frutto di un interleaving fortunato: a ogni istante
\textsc{Par-Comm} può combinare solo l'offerta corrente del ramo mittente
con quella del ramo ricevente, perché non ce n'è un terzo con cui
confondersi: non c'è scelta di interleaving da sbagliare. Questo vale
identicamente in avanti e all'indietro: invertire scambia
$\mathtt{ssend} \leftrightarrow \mathtt{srecv}$ in ciascun ramo (sopra) e ne
inverte l'ordine dei comandi ($\Inv{c_1\,;\,c_2} = \Inv{c_2}\,;\,\Inv{c_1}$),
ma non introduce mai un terzo candidato partner. È il complemento
all'indietro della confluenza dimostrata in \S\ref{sec:par-conf} per gli
interleaving forward che non toccano canali: lì la confluenza viene dalla
\emph{disgiunzione} (rami che non toccano canali operano su porzioni
disgiunte dello store, i loro passi commutano); qui viene dalla
\textbf{linearità} del canale (esattamente due estremità, \S\ref{sec:chan}),
che rende l'unico partner possibile sempre univoco. In entrambi i casi non
serve registrare la storia dell'esecuzione: la doppia guardia la evita nei
casi sequenziali di \texttt{if}/\texttt{from} (sopra), la disgiunzione delle
variabili la evita nel \texttt{par} forward, e la sessione binaria la evita
nella sincronizzazione di canale, in avanti e all'indietro.

\medskip
\noindent
Il costrutto \texttt{try\,\dots\,rollback} (\S\ref{sec:try}) non è un comando
reversibile primitivo: è una \emph{struttura di controllo} che \emph{usa}
l'inversione (esegue $\Inv{c}$ per annullare il tentativo fallito). Essendo
zucchero per \texttt{if} + \texttt{call}/\texttt{uncall}, il suo inverso eredita
quello del condizionale con chiamate.

% ======================================================================
\section{Quadro riassuntivo delle regole}
% ======================================================================

\begin{table}[H]
\centering
\small
\renewcommand{\arraystretch}{1.3}
\begin{tabular}{@{}>{$}l<{$}p{4.6cm}>{$}l<{$}@{}}
\hline
\multicolumn{1}{@{}l}{\textbf{Comando}} & \textbf{Effetto sullo store} &
\multicolumn{1}{l@{}}{\textbf{Inverso}}\\
\hline
v \mathbin{+\!\!=} e & $\sigma(v) \mathrel{+}= \mathit{eval}(e)$ & v \mathbin{-\!\!=} e\\
v \mathbin{-\!\!=} e & $\sigma(v) \mathrel{-}= \mathit{eval}(e)$ & v \mathbin{+\!\!=} e\\
v \mathbin{\hat{}\!=} e & $\sigma(v) \mathrel{\oplus}= \mathit{eval}(e)$ & \text{sé stesso}\\
v_1 \mathbin{<\!=\!>} v_2 & scambia $v_1, v_2$ & \text{sé stesso}\\
\mathtt{local}\ v = m & alloca $v = m$ & \mathtt{delocal}\ v = m\\
\mathtt{delocal}\ v = m & verifica $v = m$, rimuove & \mathtt{local}\ v = m\\
\mathtt{if}\ b \dots \mathtt{fi}\ b' & ramo per $b$, asserisce $b'$ & \mathtt{if}\ b' \dots \mathtt{fi}\ b\\
\begin{array}[t]{@{}l@{}}\mathtt{from}\ b_1\ \mathtt{do}\ c_1\\ \quad\mathtt{loop}\ c_2\ \mathtt{until}\ b_2\end{array} & ciclo, traccia $c_1[c_2\,c_1]^{*}$ & \begin{array}[t]{@{}l@{}}\mathtt{from}\ b_2\ \mathtt{do}\ \Inv{c_1}\\ \quad\mathtt{loop}\ \Inv{c_2}\ \mathtt{until}\ b_1\end{array}\\
\mathtt{call}\ f & esegue $c_f$ & \mathtt{uncall}\ f\\
\mathtt{push}(v,s) & cima $s$, azzera $v$ & \mathtt{pop}(v,s)\\
\mathtt{pop}(v,s) & verifica $v=0$, estrae cima in $v$ & \mathtt{push}(v,s)\\
\mathtt{ssend}(\langle \bar v \rangle, \mathit{ch}) & (con partner pronto) azzera $\bar v$ & \mathtt{srecv}(\langle \bar v \rangle, \mathit{ch})\\
\mathtt{srecv}(\langle \bar w \rangle, \mathit{ch}) & (con partner pronto) verifica $\bar w=\bar 0$, riceve & \mathtt{ssend}(\langle \bar w \rangle, \mathit{ch})\\
\mathtt{par}\ c_1\ \mathtt{and}\ c_2\ \mathtt{rap} & interleaving di $c_1,c_2$ (non det.) & \mathtt{par}\ \Inv{c_1}\ \mathtt{and}\ \Inv{c_2}\ \mathtt{rap}\\
\mathtt{try}\ b\ c\ \mathtt{rollback}\ r\ \mathtt{yrt}\ b' & commit se $b'$, altrimenti $\Inv{c}$ poi $r$ & \text{zucchero (\S\ref{sec:try})}\\
\hline
\end{tabular}
\caption{Le regole di riduzione di Kairos, raggruppate per costrutto.}
\label{tab:regole}
\end{table}

\noindent
Le righe $\mathtt{ssend}$/$\mathtt{srecv}$ non riducono mai da sole: l'effetto
indicato si realizza solo insieme, nella sincronizzazione congiunta
\textsc{Par-Comm} (\S\ref{sec:chan}, \S\ref{sec:par}).

\medskip
\noindent
Ogni fallimento di premessa produce $\err$ (stato locale) o blocco (stato
condiviso, canale):

\begin{table}[H]
\centering
\renewcommand{\arraystretch}{1.25}
\begin{tabular}{l l l}
\hline
\textbf{Condizione di fallimento} & \textbf{Regola} & \textbf{Esito}\\
\hline
$v \in \mathrm{dom}(\sigma)$ (\texttt{local}) & \textcolor{red!60!black}{\textsc{Local-Err}} & $\err$\\
$v \notin \mathrm{dom}(\sigma)$ o $\sigma(v) \neq m$ (\texttt{delocal}) & \textcolor{red!60!black}{\textsc{Deloc-Err}} & $\err$\\
test vero, asserzione falsa (\texttt{if}) & \textcolor{red!60!black}{\textsc{IfError1}} & $\err$\\
test falso, asserzione vera (\texttt{if}) & \textcolor{red!60!black}{\textsc{IfError2}} & $\err$\\
$\sigma(v) = 0 \land \sigma(s) = \varepsilon$ (\texttt{pop}) & \textcolor{red!60!black}{\textsc{Pop-Err}} & $\err$\\
$\sigma(v) \neq 0$ (\texttt{pop}) & \textcolor{red!60!black}{\textsc{Pop-Err2}} & $\err$\\
$\sigma(\bar w) \neq \bar 0$ (\texttt{srecv}) & \textcolor{red!60!black}{\textsc{Srecv-Err}} & $\err$\\
nessun ramo partner pronto (\texttt{ssend}/\texttt{srecv}) & $-$ & blocco\\
$b_1$ falsa all'ingresso (\texttt{from}) & \textcolor{red!60!black}{\textsc{LoopError1}} & $\err$\\
$b_1$ vera a metà ciclo & \textcolor{red!60!black}{\textsc{LoopError2}} & $\err$\\
indice fuori dai limiti ($a[k]$) & \textcolor{red!60!black}{\textsc{Idx-Err}} & $\err$\\
divisore nullo ($e/e'$, $e\,\%\,e'$) & \textcolor{red!60!black}{\textsc{Div-Err}} & $\err$\\
\hline
\end{tabular}
\caption{Condizioni di fallimento e loro esito: errore o blocco.}
\label{tab:fallimenti}
\end{table}

% ======================================================================
\paragraph{Array.} L'inversione non distingue una variabile da una cella: il
bersaglio è un luogo, e il luogo resta lo stesso.
\[
\Inv{a[e_i] \mathbin{+\!\!=} e} = a[e_i] \mathbin{-\!\!=} e
\qquad
\Inv{a[e_1] \mathbin{<\!=\!>} a[e_2]} = a[e_1] \mathbin{<\!=\!>} a[e_2]
\]
\[
\Inv{\mathtt{local}\ a[n] = m} = \mathtt{delocal}\ a[n] = m
\qquad
\Inv{\mathtt{delocal}\ a[n] = m} = \mathtt{local}\ a[n] = m
\]
L'indice $e_i$ non viene invertito perché non è un comando: è un'espressione,
valutata nello stesso store in entrambe le direzioni. Vale però la condizione
che rende l'inversione corretta: fra il comando diretto e il suo inverso le
variabili che compaiono in $e_i$ non devono essere cambiate, ed è la stessa
condizione, già richiesta per l'espressione a destra, che il vincolo statico di
\S\ref{sec:assign} garantisce.

\section{\texttt{show} e output osservabile}
\label{sec:show}
% ======================================================================

Kairos dispone di un costrutto $\mathtt{show}(v)$ che stampa il valore di una
variabile su un flusso esterno. \textbf{Non} fa parte del nucleo semantico della
macchina reversibile e per questo è stato omesso dalle regole precedenti: serve
come strumento di \emph{osservazione esterna} (ispezione dei valori durante
l'esecuzione). La transizione associata è semplicemente
$\cfg{\mathtt{show}(v)}{\sigma} \toc \cfg{\mathtt{skip}}{\sigma}$: lo store non viene toccato,
quindi $\mathtt{show}$ è trasparente rispetto a tutte le proprietà enunciate
nelle sezioni precedenti, che parlano solo di store.

\medskip
\noindent
Resta però la domanda su cosa succeda all'\textbf{output osservabile} quando
si esegue una procedura all'indietro. La risposta non è che l'inverso di
$\mathtt{show}$ sia l'identità: è che $\mathtt{show}$ è un'\emph{involuzione},
\[
\Inv{\mathtt{show}(v)} = \mathtt{show}(v),
\]
sintatticamente autoinversa come $\mathbin{\hat{}\!=}$ e $\mathbin{<\!=\!>}$, con la differenza che per quelle l'involuzione è anche semantica, mentre $\mathtt{show}\,;\,\mathtt{show}$ stampa due volte
(\S\ref{sec:rev}). Un $\mathtt{uncall}$ non sopprime dunque le stampe: le
riesegue.

\medskip
\noindent
\textbf{Nel frammento sequenziale} (senza $\mathtt{par}$), dove $\toc$ è
deterministica (\S\ref{sec:rev}), l'inversione rovescia l'ordine dei comandi
($\Inv{c_1\,;\,c_2} = \Inv{c_2}\,;\,\Inv{c_1}$), quindi i $\mathtt{show}$
vengono incontrati in ordine opposto: se in avanti la procedura stampa
$n_1, n_2, \dots, n_k$, all'indietro stampa $n_k, \dots, n_2, n_1$. È una
\emph{simmetria osservabile}, la traccia esterna che rende visibile la
reversibilità.

\medskip
\noindent
\textbf{Con $\mathtt{par}$ la simmetria è solo locale.} L'inversione agisce
ramo per ramo,
\[
\Inv{\mathtt{par}\ c_1\ \mathtt{and}\ c_2\ \mathtt{rap}} =
\mathtt{par}\ \Inv{c_1}\ \mathtt{and}\ \Inv{c_2}\ \mathtt{rap}
\quad (\S\ref{sec:rev}),
\]
quindi la sequenza di stampe \emph{di ciascun ramo} è
rovesciata; ma l'ordine con cui le stampe dei rami diversi si intercalano non
è determinato, perché l'interleaving non lo è (\S\ref{sec:par}). Già
$\mathtt{par}\ \mathtt{show}(a)\ \mathtt{and}\ \mathtt{show}(b)\ \mathtt{rap}$
può emettere $a,b$ sia in avanti sia all'indietro. L'output globale non è
quindi il rovesciamento esatto di quello forward: lo è solo la proiezione su
ogni singolo ramo.

\medskip
\noindent
Resta vero, come detto sopra, che $\mathtt{show}$ non modifica lo store e non
incide su nessuna delle proprietà enunciate, che parlano solo di store.

\section{Conformità a Janus}
\label{sec:conformita}

Che Kairos contenga Janus è un'affermazione verificabile, e conviene verificarla
su programmi che non abbiamo scritto noi. Il report originale del
linguaggio~\cite{LutzDerby1982} ne contiene tre: la radice intera, un
ordinamento a bolle con permutazione e la fattorizzazione di un intero. Sono del
1982, e la loro ortografia non è quella di oggi — il commento è \texttt{;}, lo
scambio \texttt{:}, la disuguaglianza \texttt{\#}, il resto
\texttt{\textbackslash}, e le variabili sono globali invece che parametri — ma
la struttura è la stessa, e tradurli è una sostituzione di simboli, non una
riscrittura.

\paragraph{La radice intera gira in entrambe le direzioni.} Consuma
$\mathit{num}$ producendo $\mathit{rt} = \lfloor\sqrt{\mathit{num}}\rfloor$; la
$\mathtt{uncall}$ riporta $\mathit{rt}$ a zero e $\mathit{num}$ al valore
iniziale. Da sola esercita buona parte di ciò che serve a un programma Janus:
moltiplicazione, divisione e resto dentro le guardie (\S\ref{sec:sintassi}),
cicli con la sola clausola $\mathtt{loop}$, e una procedura ausiliaria invocata
con $\mathtt{call}$ e $\mathtt{uncall}$ dentro lo stesso corpo.

\paragraph{L'ordinamento a bolle non è invertibile.} È il caso interessante,
perché il difetto sta nel listato del 1982 e non nella traduzione. Il suo
condizionale è

\begin{lstlisting}
if   list[j] > list[j+1]
then list[j] <=> list[j+1]
     perm[j] <=> perm[j+1]
fi   perm[j] > perm[j+1]
\end{lstlisting}

\noindent
e la guardia d'uscita deve distinguere i due rami: vera dopo $\mathtt{then}$,
falsa dopo $\mathtt{else}$ (\S\ref{sec:if}). Non lo fa. Con
$\mathit{list} = [4,2,5,1,3]$ e $\mathit{perm}$ identità, alla terza iterazione
esterna si ha $\mathit{list} = [1,2,4,3,5]$ e $\mathit{perm} = [3,1,0,4,2]$: il
confronto $\mathit{list}[3] > \mathit{list}[4]$ è $3 > 5$, falso, quindi si
prende il ramo vuoto, ma $\mathit{perm}[3] > \mathit{perm}[4]$ è $4 > 2$, vero.
Il condizionale non ha inverso in quel punto, perché $\mathit{perm}$ registra da
dove viene ogni elemento e non se lo scambio è appena avvenuto, e le due cose
non coincidono. È il caso d'uso dell'asserzione: senza, il programma
proseguirebbe e l'errore comparirebbe molto più tardi, come un valore sbagliato
alla chiusura di una variabile. Del listato resta perciò la sola procedura che
costruisce la permutazione identità, che è un programma Janus valido e usa un
array come parametro: la chiamata lo riempie con $0,\dots,n-1$, la
$\mathtt{uncall}$ lo riazzera. La fattorizzazione, il terzo programma, non è
stata tradotta.

\medskip
\noindent
Ai due si aggiunge un programma scritto per l'occasione, che esercita tutte e
sei le forme delle clausole facoltative di $\mathtt{if}$ e $\mathtt{from}$
(\S\ref{sec:sintassi}) con le asserzioni che devono reggere. I tre stanno nel
corpus di verifica insieme agli altri, quindi la conformità non è
un'osservazione fatta una volta ma un controllo che accompagna ogni modifica.

% ======================================================================
\chapter{Programmare in Kairos: compressione lossless reversibile}
\label{chap:programmare}
\label{sec:prog}
% ======================================================================

\section{Il punto di partenza}

La compressione \emph{lossless} è il banco di prova naturale per un linguaggio
reversibile: comprimere e decomprimere sono per costruzione l'una l'inversa
dell'altra, quindi un unico programma reversibile le realizza entrambe e non
esiste il rischio che le due metà divergano.

Il lavoro da cui partiamo è quello di Lyngby et al.~\cite{LyngbyEtAl2024}, che
ha realizzato in Janus versioni reversibili di due algoritmi dello standard zip:
la compressione a dizionario LZW e la trasformata di Burrows--Wheeler. Il loro
obiettivo non era la sola reversibilità ma la proprietà più forte di essere
\emph{clean}, cioè non lasciare spazzatura: un algoritmo reversibile non può
cancellare, quindi le strutture ausiliarie vanno azzerate invece che buttate. La
tecnica con cui ci riescono è una coppia di procedure: la prima produce
l'uscita insieme ai dati ausiliari, la seconda ricalcola gli stessi ausiliari
partendo dall'uscita, e chiamare la prima disfacendo la seconda lascia solo
l'uscita, senza storia globale.

I risultati a cui arrivano sono netti. Il loro LZW gira in $\Theta(n)$ come la
migliore versione irreversibile, quindi è asintoticamente ottimo. La loro BWT,
con l'ordinamento ingenuo delle rotazioni, costa $O(n^3)$. Concludono che le
tecniche della programmazione reversibile reggono su problemi reali, e che
restano miglioramenti \emph{possibili e necessari}, indicando in particolare
l'ordinamento delle rotazioni.

Il loro lavoro è però interamente \textbf{sequenziale}, e nel descrivere la BWT
osservano che in pratica l'ingresso si divide in blocchi, ciascuno trasformato
separatamente. È l'osservazione da cui partiamo: blocchi trasformati
separatamente sono blocchi indipendenti, e blocchi indipendenti sono rami di un
\texttt{par} che operano su porzioni disgiunte dello store, cioè l'ipotesi sotto
cui vale la confluenza (\cref{sec:par-conf}). In Janus quella strada non si può
percorrere, perché il linguaggio non ha concorrenza. In Kairos sì, ed è il
contributo di questo capitolo.

\section{Che cosa cambia scrivendolo in Kairos}

Gli algoritmi di compressione lavorano per indice, e l'accesso per indice è
diretto: \texttt{blocco[j]} legge la cella $j$ senza consumarla e senza toccare
le altre, come nel lavoro di riferimento. Il ciclo che confronta due rotazioni
legge quindi il blocco al costo di una lettura, non di una passeggiata.

\begin{lstlisting}
procedure car(int blocco[], int n, int off, int k, int c)
    local int j = 0
        j += off
        j += k
        if j >= n then
            j -= n
        fi j < off
        c += blocco[j]
        if j < off then
            j += n
        fi j >= n
        j -= k
        j -= off
    delocal int j = 0
\end{lstlisting}

\noindent
Il carattere $k$-esimo della rotazione che parte da $\mathit{off}$ sta in
posizione $(\mathit{off}+k) \bmod n$. Il modulo non si scrive con l'operatore
ma con un $\mathtt{if}$, e la ragione è la reversibilità: $j$ è una variabile
locale che va riportata a zero, quindi la correzione $j \mathrel{-}= n$ deve
poter essere disfatta, e per disfarla serve sapere se è avvenuta. La guardia
d'uscita $j < \mathit{off}$ lo dice — dopo la correzione $j$ è sceso sotto
$\mathit{off}$, senza correzione no — e il secondo $\mathtt{if}$, che è
l'inverso del primo, la usa per rimettere $j$ a posto.

\section{Il confronto lessicografico}

È il punto che il lavoro di riferimento indica come la difficoltà principale, e
lo si tocca con mano. Confrontare due rotazioni significa trovare il primo
carattere in cui differiscono, ma \textbf{non si può uscire in anticipo}: la
guardia d'uscita di un $\mathtt{if}$ non distinguerebbe ``sono entrato e non ho
ancora trovato'' da ``non sono entrato''.

La via che funziona usa un testimone. Finché le due rotazioni coincidono, il
contatore \texttt{pref} vale esattamente $k$; dopo il corpo vale $k$ oppure
$k+1$, in entrambi i casi $\geq k$, mentre chi non è entrato è rimasto sotto.

\begin{lstlisting}
if pref == k then            // siamo ancora nel prefisso comune
    if c1 > c2 then
        gt += 1
    fi gt == 1
fi pref >= k                 // guardia d'uscita valida
\end{lstlisting}

\noindent
Una formulazione alternativa, accumulare il valore posizionale in base $B$ della
differenza fra le rotazioni, ordina anch'essa correttamente ed è pura
accumulazione, ma richiede di moltiplicare per $B$ a ogni passo: in Kairos
azzerare il fattore vorrebbe una divisione, cioè sottrazioni ripetute il cui
costo dipende dal valore e non dalla lunghezza, e su numeri dell'ordine di $B^n$
non è praticabile.

\section{I blocchi in parallelo}

Ogni blocco è un ramo, con le proprie strutture; l'unico punto di contatto è il
canale su cui il ramo consegna il risultato. Il payload di un $\mathtt{ssend}$
può essere uno stack intero, quindi la colonna trasformata viaggia in un colpo
solo insieme al suo indice.

\begin{lstlisting}
procedure lavoratore(stack blocco, stack aux, stack offs, stack scarto,
                     stack out, channel c, int n)
    local int idx = 0
        call bwt(blocco, aux, offs, scarto, out, n, idx)
        ssend(<idx, out>, c)
    delocal int idx = 0

par
    call lavoratore(b0, aux0, offs0, sc0, out0, c0, n)
and
    call lavoratore(b1, aux1, offs1, sc1, out1, c1, n)
and
    call raccoglitore(c0, c1, col0, col1, indici)
rap
\end{lstlisting}

\noindent
Qui il blocco viaggia come \texttt{stack} e non come array, e la ragione è il
canale: il payload di un $\mathtt{ssend}$ può essere un intero, uno stack o un
canale, non un vettore. Trasferire un array richiederebbe di deciderne la sorte
nel mittente — un rendez-vous consuma ciò che manda, e un array ha lunghezza
fissa — ed è una scelta che non abbiamo compiuto. La versione concorrente
resta perciò quella su stack, e i tempi che seguono sono i suoi.

\section{Risultati}

\paragraph{Correttezza.} Sul vettore di prova del lavoro di riferimento, la
stringa \texttt{banana} con alfabeto ridotto ($a=1$, $b=2$, $n=3$):

\begin{center}
\begin{tabular}{ll}
\hline
blocco & \texttt{\{2, 1, 3, 1, 3, 1\}}\\
rotazioni ordinate & \texttt{\{5, 3, 1, 0, 4, 2\}}\\
ultima colonna & \texttt{\{3, 3, 2, 1, 1, 1\}} \quad cioè \texttt{nnbaaa}\\
indice & \texttt{3}\\
ricostruzione da (colonna, indice) & \texttt{\{1, 3, 1, 3, 1, 2\}}\\
\hline
\end{tabular}
\end{center}

\noindent
Coincide con il loro esempio. La ricostruzione parte da \emph{sola} colonna e
indice, quindi è decompressione vera e non un semplice \texttt{uncall}: usa la
trasformata inversa classica, che costruisce il vettore delle righe precedenti e
lo percorre a partire dall'indice. Il cammino attraversa la matrice a partire
dalla riga $\mathit{idx}$, quindi restituisce il blocco a meno di una rotazione:
$\{1,3,1,3,1,2\}$ è \texttt{banana} letta da capo dopo la \texttt{b}. Verificata anche su blocchi di dodici simboli
e tre blocchi concorrenti.

\paragraph{Costo.} Il tempo di esecuzione di un blocco, al variare della sua
lunghezza $n$ (bytecode congelato, misura del solo interprete, minimo di cinque
ripetizioni):

\begin{center}
\begin{minipage}[c]{0.28\textwidth}
\centering
\begin{tabular}{ll}
\hline
$n$ & tempo\\
\hline
12 & 0{,}04 s\\
16 & 0{,}10 s\\
20 & 0{,}20 s\\
24 & 0{,}35 s\\
28 & 0{,}56 s\\
32 & 0{,}87 s\\
\hline
\end{tabular}
\end{minipage}%
\hfill
\begin{minipage}[c]{0.66\textwidth}
\centering
% Assi logaritmici in base 2: su di essi una legge di potenza e' una retta, e
% l'esponente e' la pendenza. La retta disegnata e' quella dei minimi quadrati.
\begin{tikzpicture}[x=3.3cm,y=1.05cm]
  \draw[->] (-0.08,0) -- (1.72,0) node[right,font=\footnotesize] {$n$ (log)};
  \draw[->] (-0.08,-0.15) -- (-0.08,5.35) node[above right,font=\footnotesize] {tempo (log)};
  % tacche su n
  \foreach \X/\L in {0.085/12, 0.5/16, 0.822/20, 1.085/24, 1.307/28, 1.5/32}
    \draw (\X,0) -- (\X,-0.09) node[below,font=\footnotesize] {\L};
  % tacche sul tempo (secondi, raddoppi)
  \foreach \Y/\L in {1/{0{,}0625\,s}, 2/{0{,}125\,s}, 3/{0{,}25\,s}, 4/{0{,}5\,s}, 5/{1\,s}}
    \draw (-0.08,\Y) -- (-0.13,\Y) node[left,font=\footnotesize] {\L};
  % retta dei minimi quadrati, pendenza 2,84
  \draw[thick,gray!70] (0,0.219) -- (1.6,5.067);
  % punti misurati
  \foreach \X/\Y in {0.085/0.500, 0.5/1.725, 0.822/2.692, 1.085/3.491, 1.307/4.168, 1.5/4.792}
    \fill (\X,\Y) circle (1.7pt);
  \node[font=\footnotesize,anchor=west] at (0.62,1.30) {pendenza $3{,}03$};
\end{tikzpicture}
\end{minipage}
\end{center}

\noindent
Gli assi sono logaritmici, quindi l'allineamento dei punti dice che il costo è
una legge di potenza e la pendenza ne dà l'esponente. La retta di regressione di
$\log t$ su $\log n$ ha pendenza $3{,}03$: il costo è $n^3$, lo stesso del lavoro
di riferimento, e l'accesso per indice non aggiunge nulla all'ordine di
grandezza. È quello che ci si aspetta dalla forma
del programma: l'ordinamento a bolle compie $O(n^2)$ confronti fra rotazioni, e
ciascun confronto ne scandisce $O(n)$ caratteri.

\paragraph{Concorrenza.} Lo stesso calcolo in due varianti che differiscono in
un solo punto, la procedura che orchestra i blocchi. Nella prima ogni blocco sta
in un \texttt{par} a sé, insieme alla sola ricezione del proprio risultato: il
lavoratore e il suo ricevente procedono insieme, ma i blocchi si susseguono, e
non c'è parallelismo \emph{fra} blocchi. Nella seconda tutti i lavoratori sono
rami di un unico \texttt{par}, con il raccoglitore come ultimo ramo. Con $n=8$,
su una macchina a \emph{due} nuclei:

\begin{center}
\begin{tabular}{llll}
\hline
blocchi & un \texttt{par} per blocco (s) & un \texttt{par} solo (s) & guadagno\\
\hline
2 & 0{,}71 & 0{,}62 & 1{,}16$\times$\\
4 & 1{,}49 & 1{,}11 & 1{,}34$\times$\\
6 & 2{,}21 & 1{,}67 & 1{,}32$\times$\\
8 & 2{,}87 & 1{,}96 & 1{,}46$\times$\\
\hline
\end{tabular}
\end{center}

\noindent
Il dato da leggere non è il guadagno in sé ma quanto si avvicina al massimo che
la macchina rende disponibile. Con due nuclei quel massimo è $2\times$, e a otto
rami se ne raggiunge il settantatré per cento: il parallelismo che il linguaggio
dichiara si realizza. Invertendo la legge di Amdahl, un guadagno di $1{,}46\times$
su due processori corrisponde a una frazione parallela del sessantatré per cento
circa, e la parte restante ha un candidato preciso nella forma del programma: il
raccoglitore è un ramo solo e riceve dai canali uno per volta, quindi le consegne
si allineano.

Va detto ciò che questi dati non dicono. Sono presi su due nuclei, e da soli non
autorizzano alcuna previsione sul comportamento con molti più nuclei: la stessa
frazione sequenziale ammetterebbe in linea di principio guadagni ben superiori a
quelli qui osservati. Il guadagno cresce con il numero di rami senza appiattirsi
entro i quattro punti misurati, e con quattro volte più rami che nuclei ci si
aspetterebbe di vedere la contesa fra thread: quattro punti non bastano a dire
se manchi o se cominci più in là. Stabilire dove il guadagno si fermi davvero richiede
la misura su una macchina con più nuclei, e un raccoglitore ad albero per togliere
di mezzo la parte sequenziale del programma di prova.

\section{Limiti}

Tre, dichiarati. L'alfabeto delle prove è ridotto a tre simboli: nulla nel
codice ne dipende, ma neppure lo esercita. I blocchi sono corti, imposti dal
costo $n^3$ dell'ordinamento ingenuo delle rotazioni, che è il punto su cui il
lavoro di riferimento stesso dichiara che i miglioramenti restano necessari. E
la proprietà \emph{clean} nel senso stretto, cioè
azzerare anche l'ingresso e non solo le strutture ausiliarie, richiederebbe la
coppia di procedure applicata all'intera trasformata: qui gli ausiliari si
azzerano disfacendo le procedure che li hanno prodotti, l'ingresso resta.

% ======================================================================
\chapter{Implementazione: interprete e DAP}
\label{chap:interprete}
% ======================================================================

\section{La catena, dal sorgente all'esecuzione}

L'implementazione è divisa in due metà, con un confine netto. La prima, in
Python, porta il sorgente fino a un bytecode testuale; la seconda, in C, esegue
quel bytecode nelle due direzioni.

\begin{center}
\begin{tabular}{lll}
\hline
& \textbf{stadio} & \textbf{prodotto}\\
\hline
\multirow{4}{*}{Python}
 & analisi lessicale   & sequenza di token\\
 & analisi sintattica  & albero sintattico\\
 & generazione         & bytecode testuale\\
\hline
C & interprete         & esecuzione, in avanti e all'indietro\\
\hline
\end{tabular}
\end{center}

\noindent
Il confine sta nel bytecode, ed è un file di testo. La scelta ha un costo, che
si vedrà in \S\ref{sec:interprete}, e due ragioni. La prima è ispezionabilità:
si può leggere ciò che la macchina esegue, e ogni riga porta il numero di riga
del sorgente da cui proviene, che è quanto serve al debugger per fermarsi nel
punto giusto. La seconda è che il testo basta a sé anche all'indietro: il corpo
inverso di una procedura è a sua volta bytecode, emesso dallo stesso generatore
e eseguito dallo stesso ciclo, e non serve una seconda macchina che lo
interpreti.

\section{Un esempio passo per passo}

Si consideri una procedura che somma i primi $n$ interi, con il ciclo a due
corpi (\cref{sec:loop}):

\begin{lstlisting}
procedure somma(int n, int s)
    local int i = 0
    from i == 0 do
        s += i
    loop
        i += 1
    until i == n
    delocal int i = n
\end{lstlisting}

\paragraph{Analisi lessicale.} Il primo stadio riduce il testo a una sequenza di
token, scartando spazi e commenti. Le parole chiave non sono espressioni
regolari a sé: si riconosce un identificatore e si guarda se compare in una
tabella di riservate, così aggiungerne una costa una riga.

\begin{lstlisting}
PROCEDURE ID(somma) LPAREN INT ID(n) COMMA INT ID(s) RPAREN
LOCAL INT ID(i) EQUALS NUMBER(0)
FROM ID(i) EQEQ NUMBER(0) DO
  ID(s) PLUSEQUALS ID(i)
LOOP
  ID(i) PLUSEQUALS NUMBER(1)
UNTIL ID(i) EQEQ ID(n)
DELOCAL INT ID(i) EQUALS ID(n)
\end{lstlisting}

\paragraph{Analisi sintattica.} La grammatica è LALR(1), e ogni produzione
costruisce direttamente il nodo corrispondente. Quella del ciclo mostra bene il
rapporto fra le due cose:

\begin{lstlisting}
statement : FROM condition DO opt_body LOOP opt_body UNTIL condition

p[0] = ('from', p[2], p[4], p[8], p.lineno(1), p.lineno(7), p[6])
//               b1    c1    b2                             c2
\end{lstlisting}

\noindent
I due numeri di riga, quella del \code{from} e quella dell'\code{until},
servono al debugger per fermarsi sull'uno o sull'altro estremo del ciclo.

\paragraph{Albero sintattico.} I nodi sono tuple, non oggetti, e ogni nodo porta
in coda il numero di riga.

\begin{lstlisting}
('procedure', 'somma', [('int','n'), ('int','s')],
 [('local', 'int', 'i', 0, 2),
  ('from', ('cond','==','i',0),          // guardia d'ingresso
           [('assign','s','+=','i',4)],  // corpo do
           ('cond','==','i','n'),        // guardia d'uscita
           3, 7,
           [('assign','i','+=',1,6)]),   // corpo loop
  ('delocal', 'int', 'i', 'n', 8)], 1)
\end{lstlisting}

\noindent
Il secondo corpo sta in coda alla tupla, dopo i numeri di riga: le procedure che
percorrono l'albero e leggono solo il primo corpo continuano a funzionare, e
vanno estese soltanto dove il secondo conta davvero.

\paragraph{Bytecode.} Ogni riga è un numero di riga sorgente, un opcode e i suoi
argomenti.

\begin{lstlisting}
@1   PROC somma
@1   PARAM int n
@1   PARAM int s
@2   LOCAL int i 0
@3   EVAL i == 0             // guardia d'ingresso
@3   ASSERTT                 // dev'essere vera qui  (LoopError1)
@3   JMPF FROM_ERR_5
@3   JMP FROM_START_5        // entra saltando il secondo corpo
@3   LABEL FROM_BACK_5
@6   PUSHEQ i 1              //   corpo loop
@3   EVAL i == 0
@3   ASSERTZ                 // e falsa qui         (LoopError2)
@3   LABEL FROM_START_5
@4   PUSHEQ s i              //   corpo do
@7   EVAL i == n             // guardia d'uscita
@7   JMPF FROM_BACK_5        // falsa: altra iterazione, passando dal loop
@7   LABEL FROM_END_5
@7   LABEL FROM_ERR_5
@8   DELOCAL int i n
@1   END_PROC somma
\end{lstlisting}

\noindent
Il salto iniziale è ciò che realizza la traccia $c_1\,[c_2\,c_1]^{*}$: alla
prima entrata il secondo corpo viene scavalcato, e ci si torna solo dal
\emph{back-edge}.

Le due asserzioni del ciclo (\cref{sec:loop}) — $b_1$ vera alla prima entrata,
\textsc{LoopError1}, e falsa a ogni ri-entrata, \textsc{LoopError2} — sono
emesse come opcode dedicati, e abortiscono l'esecuzione quando cadono. Non sono
diagnostica accessoria: è quell'invariante a rendere il ciclo invertibile per
sola sintassi, perché senza di esso l'inversa
$\mathtt{from}\ b_2\ \mathtt{do}\ \Inv{c_1}\ \mathtt{loop}\ \Inv{c_2}\
\mathtt{until}\ b_1$ uscirebbe dopo una iterazione mentre l'originale ne compie
$n$. L'implementazione di riferimento di Janus adotta la stessa scelta:
\texttt{assertTrue} prima della prima iterazione e \texttt{assertFalse} dopo
ogni corpo \emph{loop}~\cite{YokoyamaGluck2007}.

L'inverso di una procedura si ottiene in un modo solo, ed è quello della
\cref{sec:call}: si applica $\Inv{\cdot}$ all'albero sintattico e si emette un
secondo corpo, dentro la stessa procedura, che l'esecuzione diretta scavalca con
un salto e a cui $\mathtt{uncall}$ entra per etichetta. Da lì in poi
$\mathtt{uncall}$ è una normale esecuzione in avanti: stesso ciclo, stesso
frame, stessi parametri, e nessuna macchina separata che percorra il programma
al contrario cercando di dedurne la struttura.

Il $\mathtt{par}$ non fa eccezione: il corpo inverso di una procedura che ne
contiene uno è $\mathtt{par}\ \Inv{c_1}\ \mathtt{and}\ \dots\ \mathtt{and}\
\Inv{c_n}\ \mathtt{rap}$, cioè la regola del \cref{chap:reversibilita} applicata
alla lettera --- rovesciando l'ordine dentro ogni ramo si rovescia la successione
degli appuntamenti, e i due estremi di ciascuno restano l'uno di fronte
all'altro. Vale a una condizione, ed è la stessa da cui viene l'unicità
dell'accoppiamento: che i titolari di ogni canale siano fissati dal testo. Se un
canale viaggia \emph{dentro} un messaggio --- $\mathtt{ssend}(\langle x,
\mathit{ch}_2\rangle, \mathit{ch}_1)$, la mobilità del $\pi$-calcolo --- la
titolarità di $\mathit{ch}_2$ nasce invece da un appuntamento precedente su
$\mathit{ch}_1$, e rovesciando l'ordine dei comandi gli usi di $\mathit{ch}_2$
verrebbero prima dell'appuntamento che lo consegna. Non è un timore teorico: dei
programmi di prova due lo fanno, e con il corpo inverso emesso comunque si
fermano entrambi in stallo, ogni ramo in attesa su un canale di cui nessun altro
è ancora titolare. Il compilatore riconosce quindi il caso, e lo riconosce sul
programma intero perché un ramo può limitarsi a chiamare la procedura che
trasmette; per quei due programmi su ottantanove non emette corpo inverso, e
l'inversione resta affidata al motore dinamico, che accoppia gli appuntamenti
mentre esegue invece di prevederli. Per tutti gli altri, $\mathtt{par}$ inclusi,
l'inverso è compilato.

Sta dentro la stessa procedura, e non in una nuova, perché così ne condivide la
tabella dei nomi: l'ordine in cui gli indici vengono assegnati, che il dump
finale rende osservabile, resta quello dell'esecuzione diretta.

\section{Strutture e ciclo di esecuzione}
\label{sec:interprete}

\paragraph{Strutture.} Lo store $\sigma$ della semantica è un \code{Frame}: un
vettore di variabili, ciascuna con tipo, valore e, se è uno stack o un canale, la
struttura associata. Una configurazione $\langle \sigma, c\rangle$ è la coppia
frame corrente e puntatore alla riga di bytecode.

La tabella che traduce i nomi in indici non sta nel frame ma nella
\emph{procedura}, ed è condivisa da tutte le sue attivazioni. La distinzione
conta: un indice denota allora lo stesso nome ovunque, mentre una tabella per
attivazione lo legherebbe al momento in cui l'attivazione è nata, e lo stesso
indice potrebbe indicare variabili diverse in due istanze della stessa
procedura.

La traduzione nome-indice è una scansione lineare, con confronto sul primo
carattere prima di quello completo. Il numero di nomi di una procedura è
dell'ordine delle decine, e una ricerca ne tocca in media cinque, su poche linee
di cache che restano calde fra un'istruzione e la successiva: a
questa scala la scansione costa meno di un indice associativo, che andrebbe
costruito e mantenuto e i cui accessi sono sparsi.

La risoluzione non viene però ripetuta. L'indice di un nome è una costante del
programma — la tabella appartiene alla procedura ed è a sola aggiunta, quindi un
indice assegnato non cambia — e viene perciò memorizzato sull'istruzione che lo
usa: è la tecnica che gli interpreti chiamano
\emph{quickening}~\cite{BrunthalerQuickening,PEP659}, e a differenza di un
indice associativo non costruisce alcuna struttura. Vale un vincolo, e non è
formale: risolvere un nome lo \emph{inserisce}, e l'ordine di inserimento è
osservabile nell'ordine del dump finale e nel conteggio delle celle. La cache si
riempie dunque \emph{dopo} la risoluzione, nel punto in cui l'interprete la
compiva comunque, e accoglie solo i risultati positivi: un nome assente in questo
istante può comparire più tardi per una dichiarazione locale in un'altra
attivazione, e memorizzare un fallimento darebbe una risposta stantia. Un operando
che comincia per cifra non è un nome — lo garantisce il lessico — e per esso non
si esegue alcuna ricerca. Sul programma di compressione resta una ricerca di
nome ogni quattro istruzioni e mezza eseguite: le altre trovano l'indice già
scritto accanto a sé e non cercano affatto.

Nessuna struttura ha una dimensione fissata in anticipo: non la profondità di
ricorsione, non il numero di rami di un \texttt{par}, non i parametri di una
procedura, non i nomi che un frame può contenere né la lunghezza di ciascuno.
Tutte crescono su richiesta, e l'unico limite che resta, sulla lunghezza di una
riga di bytecode, è verificato e produce un errore esplicito invece di troncare.
Un frame è così duecento byte: il vettore delle variabili e i riferimenti alla
tabella dei nomi della sua procedura. Di frame ne esiste uno per procedura, non
uno per attivazione: la ricorsione non ne duplica la struttura. Un nome aperto da
$\mathtt{local}$ in un'attivazione annidata prende lo slot dell'omonimo
dell'attivazione esterna e lo tiene da parte; il $\mathtt{delocal}$ che lo chiude
rimette al suo posto quello esterno. È l'ombreggiamento della semantica
(\S\ref{sec:scope}), letto direttamente nello store invece di essere imitato con
una copia. Ciò che distingue un'attivazione dall'altra è allora solo quel poco
che la macchina già salva sulla pila delle chiamate --- i parametri legati e la
pila dei nomi aperti --- e per questo un secondo $\mathtt{local}$ sullo stesso
nome nella \emph{stessa} attivazione resta l'errore che la semantica prescrive,
mentre fra attivazioni diverse è ombreggiamento.

\paragraph{Ciclo principale.} Il testo del bytecode non cambia mai durante
l'esecuzione: è una sola istanza condivisa da tutti i rami, e una variante di
compilazione che la mappa in sola lettura lo verifica. Tutto ciò che se ne
ricava è dunque una costante, e va calcolato una volta. L'interprete lavora
perciò su un vettore di istruzioni \emph{decodificate}, in cui l'indice di
un'istruzione è il suo numero di riga meno uno: passare da un numero di riga a
un'istruzione è dunque una sottrazione, non una scansione. La decodifica è
pigra per procedura — una procedura mai chiamata non viene mai decodificata — e
le posizioni degli operandi sono scostamenti, non puntatori, perché i rami
paralleli possono lavorare su copie del programma.

\begin{lstlisting}
const Insn *in = &g_ir[ip];   /* indice = riga - 1 */
uint8_t op_tag = in->op;      /* classificato una volta sola */
...
if      (op_tag == INVOP_PUSHEQ) op_pusheq_ir(vm, cur_fi, in, ptr);
else if (op_tag == INVOP_MINEQ)  op_mineq_ir (vm, cur_fi, in, ptr);
else if (op_tag == INVOP_XOREQ)  op_xoreq_ir (vm, cur_fi, in, ptr);
\end{lstlisting}

\noindent
La classificazione dal testo resta come ripiego, per le righe che l'array non
copre; è uno \texttt{switch} sul primo carattere dell'opcode seguito dai soli
confronti che quel carattere lascia possibili, quindi il suo costo non cresce
col numero di opcode del linguaggio.

\noindent
L'inversione non richiede alcuna storia dei dati, ed è il punto in cui la
reversibilità del linguaggio si paga da sola: l'inverso di un comando è
\emph{un altro comando}, e $\Inv{\cdot}$ lo calcola dall'albero sintattico. Il
corpo inverso è quindi bytecode come l'altro, prodotto una volta in
compilazione, ed eseguito dallo stesso ciclo (\S\ref{sec:call}). Nessuna
struttura di controllo va ricostruita a tempo di esecuzione: dove comincia un
ciclo e quale ramo è stato preso sono domande che il corpo inverso non pone,
perché la risposta è già scritta nella sua forma. È l'invariante del ciclo a
renderlo possibile, ed è la ragione per cui le asserzioni non sono
diagnostica accessoria.

\section{Il parallelismo}

Un blocco \texttt{par} genera un thread di sistema per ramo. I rami non
condividono memoria per costruzione (\cref{sec:par}), quindi non serve alcun
lock sullo store; l'unica primitiva di sincronizzazione è il canale, realizzato
con un mutex e due code di attesa, una per i mittenti e una per i riceventi. Chi
arriva per primo si accoda e si blocca, chi arriva per secondo trova il partner,
lo sveglia, e lo scambio avviene: è il rendez-vous di \textsc{Par-Comm}.

\medskip
\noindent
Non condividere lo store non basta però a rendere i rami indipendenti: lo
diventano solo se non condividono nulla \emph{nemmeno nell'esecutore}. Quattro
strutture stanno sul cammino di ogni ramo, e ciascuna è per thread o si tocca
con operazioni atomiche: lo stato della tokenizzazione; il rilevatore di
scritture concorrenti su una cella intera, che vive nella cella e non in un lock
unico; l'indice dei nomi, che si consulta senza esclusione perché è a sola
aggiunta e pubblica ogni voce solo quando è completa, e prende il lock soltanto
per inserire; e la modalità di soppressione dell'output durante il replay in
avanti, che è una proprietà dell'esecuzione di un singolo ramo.

Il criterio che le accomuna non è che tocchino i dati, visto che l'ultima non li
tocca, ma che una copia sola condivisa fra i rami renderebbe il risultato
dipendente dall'interleaving. Per lo store lo vieta il linguaggio; per le strutture di
servizio dell'esecutore deve provvedere l'implementazione, e la confluenza vale
sull'osservabile solo se provvede a tutte.

\medskip
\noindent
Del bytecode esiste una sola istanza, che tutti i rami leggono. È possibile
perché il testo del programma porta soltanto istruzioni: ogni stato vive nei
frame, quindi nessuno ha ragione di scrivere nel programma, e ciò che serve a
delimitare una riga (dove comincia e quanto è lunga) sta nell'indice costruito
dopo la normalizzazione. Chi deve tokenizzare copia la riga in un buffer
proprio, perché la tokenizzazione della libreria C sostituisce i separatori con
terminatori, e quella copia è dell'ordine delle decine di byte invece che
dell'intero programma.

La proprietà è verificata e non solo dichiarata: una modalità di collaudo mette
il programma normalizzato in una regione di memoria protetta in sola lettura, e
un'eventuale scrittura fa terminare il processo nel punto esatto in cui avviene.
Sull'intero insieme dei programmi di prova non se ne verifica nessuna.

\section{Il debugger DAP}

L'interprete espone un'interfaccia di debug, e sopra di essa un'estensione per
Visual Studio Code\footnote{\url{https://github.com/nicologiuliani6/kairos-vscode-debugger}}
parla il Debug Adapter Protocol, il protocollo con cui l'editor dialoga con un
debugger qualunque. Si mettono punti di interruzione sul sorgente Kairos, si
percorre l'esecuzione in avanti e all'indietro un passo per volta o fino al
punto di interruzione successivo, si ispezionano le variabili mentre cambiano e
si legge l'output del programma nella console di debug.

Il collegamento fra i due lati è ciò che il bytecode testuale rende immediato:
ogni riga porta il numero della riga sorgente da cui proviene
(\S\ref{sec:interprete}), quindi un punto di interruzione posto sul sorgente
diventa una condizione sulle righe eseguite, senza una tabella di corrispondenza
da costruire e tenere allineata. Lo stato mostrato all'editor è lo stesso store
dell'esecuzione, serializzato quando l'editor lo chiede: non una copia tenuta
aggiornata in parallelo, ma il frame corrente letto in quel momento.

\medskip
\noindent
Il movimento all'indietro è ciò che distingue questo debugger da uno ordinario,
e non è un'aggiunta all'interfaccia: un passo indietro esegue l'\emph{inversa}
dell'ultima istruzione, cioè la stessa direzione di esecuzione che il linguaggio
possiede. L'inverso di un'operazione sui dati è un'altra operazione sui dati, e
le istruzioni di controllo non toccano lo store, quindi all'indietro sono solo
posizione. Un debugger a viaggio nel tempo in un linguaggio ordinario deve
registrare e riprodurre; qui non c'è nulla da registrare, perché l'informazione
per tornare indietro è già nello stato.

Un vincolo di realizzazione merita di essere detto, perché non è accidentale: il
passo indietro avviene dentro il ciclo dell'interprete e non nello strato che
parla il protocollo. Il puntatore all'istruzione corrente vive nel primo, e
modificare lo store dal secondo lo lascerebbe fermo dov'era, con posizione e
memoria che raccontano due storie diverse. Attraversare il confine di una
chiamata è il caso che resta fuori: richiede di ricostruire il record di
attivazione, che non è parte dello store ma della macchina che lo esegue.

% ======================================================================
\chapter{Lavori correlati}
\label{chap:correlati}
\label{sec:related}
% ======================================================================

Il lavoro più vicino per obiettivi è una precedente estensione concorrente di
Janus~\cite{Sammarchi2019}, che aggiunge al linguaggio un costrutto
\texttt{fork\,\dots\,and\,\dots\,join} e uno scambio di messaggi, e li realizza
con un compilatore verso C++ che usa i thread POSIX. Il confronto è utile
proprio perché il problema di partenza è lo stesso, e ogni scelta diversa si
lega a una proprietà diversa.

La prima differenza è di metodo: là i costrutti sono definiti dal codice in cui
vengono tradotti, qui da regole di riduzione. Avere la semantica non è un
ornamento: è ciò che permette di enunciare la confluenza e la reversibilità come
proprietà del linguaggio, e non come comportamento osservato di una particolare
traduzione.

Le altre due differenze riguardano lo stato e i canali, e sono legate fra loro.
Quel lavoro mantiene deliberatamente uno \emph{stato globale unico}, per non
dover trattare allocazione e deallocazione nei rami concorrenti, e introduce le
\texttt{struct} proprio per far arrivare dati ai thread. Kairos fa l'opposto e
pretende che i rami tocchino porzioni disgiunte dello store: è una restrizione,
ma è l'ipotesi da cui discende la confluenza degli interleaving
(\S\ref{sec:par-conf}), che con memoria condivisa non sarebbe dimostrabile. Allo
stesso modo, là i canali sono \emph{porte}, cioè code di messaggi con una
capacità massima; qui la comunicazione è un rendez-vous senza buffer. Anche
questa è una rinuncia, ed è quella che rende $\mathtt{ssend}$ e $\mathtt{srecv}$
una coppia inversa esatta: su una coda i due agirebbero su estremità opposte, e
l'inverso di un accodamento in fondo non sarebbe una ricezione dalla testa
(\S\ref{sec:chan}).

Fra i linguaggi reversibili e concorrenti il più vicino per costruzione è invece
\textbf{CRIL}~\cite{OguchiYuen2024}, che però lavora a un livello di astrazione
diverso: è un \emph{linguaggio intermedio}, non un linguaggio in cui si scrive.
Un programma CRIL è una
collezione di \emph{basic block}, ciascuno costituito da una singola istruzione a
tre indirizzi con etichette per il flusso di controllo in avanti e all'indietro.
CRIL estende RIL di Mogensen permettendo l'esecuzione concorrente di più
chiamate di basic block in entrambe le direzioni, e la sua semantica operazionale
soddisfa \emph{causal safety} e \emph{causal liveness} come criteri di
correttezza per la reversibilità.

Le differenze rispetto a Kairos sono sistematiche, e riguardano il livello di
astrazione prima ancora dei costrutti:

\begin{table}[H]
\centering
\renewcommand{\arraystretch}{1.25}
\small
\begin{tabular}{l l l}
\hline
 & \textbf{CRIL}~\cite{OguchiYuen2024} & \textbf{Kairos}\\
\hline
Livello & intermedio (basic block, 3 indirizzi) & alto, strutturato (stile Janus)\\
Concorrenza & chiamate concorrenti di blocchi & \texttt{par\,\dots\,and\,\dots\,rap} esplicito\\
Comunicazione & variabili condivise & canali sincroni, store \emph{disgiunti}\\
Correttezza & causal safety / causal liveness & reversibilità (\S\ref{sec:rev}), confluenza (\S\ref{sec:par-conf})\\
\hline
\end{tabular}
\caption{Kairos e CRIL a confronto.}
\label{tab:cril}
\end{table}

\noindent
La scelta di comunicare esclusivamente per canali su porzioni disgiunte dello
store, invece che per variabili condivise, non è incidentale: è precisamente
l'ipotesi di disgiunzione da cui discende la confluenza degli interleaving che
non toccano canali (\S\ref{sec:par-conf}). Un confronto puntuale fra i due
criteri di correttezza (causal safety/liveness da un lato, la proprietà di
reversibilità di \S\ref{sec:rev} dall'altro) resta da svolgere.

\medskip
\noindent
Sulle \textbf{transazioni} il termine di paragone è invece la linea di lavori su
$\rho\pi$ e croll-$\pi$ \cite{LaneseMezzinaSchmittStefani2011,%
LaneseLienhardtMezzinaSchmittStefani2013}, dove il ritorno all'indietro è
governato da primitive esplicite e produce esattamente il comportamento di una
transazione con compensazione. La differenza sta in dove si ferma la
propagazione. In quei calcoli un rollback si propaga fra i processi: se un
processo che deve essere disfatto ha nel frattempo comunicato con un altro,
anche l'altro va disfatto, e il meccanismo che governa questa propagazione è
buona parte della loro difficoltà. Kairos evita il problema per costruzione,
vietando comunicazione e concorrenza dentro il corpo di un
$\mathtt{try}\ldots\mathtt{rollback}$ (\S\ref{sec:try}): una transazione è
locale a un ramo, e non c'è nulla da propagare. È una restrizione forte, ed è
anche la ragione per cui le nostre transazioni si riducono a zucchero sintattico
su costrutti già presenti; toglierla significherebbe affrontare la propagazione
del rollback, che è il contributo proprio di quei lavori.

% ======================================================================
\chapter{Conclusioni e lavori futuri}
\label{chap:conclusioni}
\label{sec:concl}
% ======================================================================

Abbiamo presentato Kairos, un linguaggio imperativo strutturato che è insieme
reversibile e concorrente, e ne abbiamo dato la semantica operazionale
strutturale in stile small-step, definita come estensione di quella di Janus di
Lami, Lanese e Stefani. Ciò che Kairos aggiunge alla base sono i costrutti che
servono a un linguaggio concorrente e a un linguaggio con dati: composizione
parallela, comunicazione sincrona su canali tipizzati, transazioni con
compensazione, stack, parametri di procedura e ambito locale. Abbiamo definito
l'operatore di inversione, sintattico e calcolabile in tempo lineare, e mostrato
che l'esecuzione all'indietro di un blocco parallelo è ben definita: la
disgiunzione delle porzioni di store toccate dai rami dà la confluenza degli
interleaving che non comunicano, e il vincolo di sessione binaria rende unico
l'accoppiamento fra un invio e la ricezione che lo consuma. Quel vincolo è
sorvegliato come proprietà dell'esecuzione e non del testo, il che consente la
delega, cioè il passaggio di un canale da un titolare a un altro. Il linguaggio
è realizzato da un interprete, e su di esso abbiamo condotto un caso di studio:
la trasformata di Burrows--Wheeler reversibile, con i blocchi trasformati in
parallelo, verificata sul vettore di prova del lavoro da cui parte e misurata.

\paragraph{Maturità pratica.} Il divario più netto è fra ciò che la teoria dice
possibile e ciò che servirebbe a un compressore vero, e i due algoritmi dello
standard zip lo mostrano in modo opposto. La compressione a dizionario regge:
nella versione reversibile di riferimento gira in $\Theta(n)$, come la migliore
irreversibile, quindi la reversibilità non costa nulla in complessità. La
trasformata no: costa $O(n^3)$, ed è l'esponente che misuriamo. Resta però
$O(n^3)$, cioè l'ordinamento ingenuo delle rotazioni, che è il punto su cui il
lavoro di riferimento stesso dichiara che i miglioramenti sono possibili e
necessari: un ordinamento migliore è la direzione che indicheremmo per prima.
Finché manca, i blocchi restano corti, e con blocchi corti la trasformata non ha
il materiale su cui produrre l'effetto per cui esiste. Manca inoltre la
proprietà \emph{clean} nel senso stretto: le strutture ausiliarie si azzerano
disfacendo le procedure che le hanno prodotte, ma l'ingresso sopravvive
all'uscita, mentre un compressore deve consumarlo.

\paragraph{Il linguaggio.} La disciplina di sessione è parziale per costruzione:
mancando la scelta etichettata, i due rami di un $\mathtt{if}$ che comunichi
devono compiere le stesse azioni sul canale, e dove il frammento non basta il
tipo resta indeterminato. Aggiungere la scelta renderebbe totale ciò che oggi
riconosce i soli conflitti dimostrabili; i protocolli multiparty estenderebbero
il tutto a canali con più di due partecipanti. Più in generale, l'assenza di
corse critiche è qui un'ipotesi sotto cui valgono le proprietà, non un teorema:
darle la forma di un sistema di tipi, separato dal calcolo, è il modo consueto
di ottenerla e resta da fare. La comunicazione asincrona è un'altra estensione
naturale, con l'avvertenza che non è un'aggiunta indolore: cambia ciò che si
può dimostrare sull'inversione, perché un messaggio in transito è stato che il
rendez-vous non produce.

\paragraph{L'implementazione.} Sul calcolo concorrente la misura dice che il
parallelismo si realizza, con l'ottantanove per cento del massimo che la macchina
di prova rende disponibile, e indica nella frazione sequenziale del programma, cioè
nel raccoglitore che riceve dai canali uno per volta, il candidato per la parte
restante. Non dice però fin dove il guadagno cresca: servono una macchina con più
nuclei e un raccoglitore ad albero, e sono due esperimenti, non due congetture.

Nella direzione opposta, quella in cui l'interprete aderisce al modello, ciò che
la semantica descrive come una sola istanza è una sola istanza anche
nell'esecutore. Il testo del programma è condiviso, perché porta solo
istruzioni e nessuno stato; l'inversione lo è a sua volta, perché il corpo
inverso è bytecode prodotto in compilazione e $\mathtt{uncall}$ non costruisce
alcun contesto temporaneo; e lo store è una sola istanza anche lui. Un frame
appartiene alla procedura, non all'attivazione: la ricorsione non ne duplica la
struttura, e i nomi che $\mathtt{local}$ apre in un'attivazione ombreggiano gli
omonimi di quelle esterne, come nella semantica. Ciò che resta separato è
separato perché il linguaggio lo prescrive: i rami di un $\mathtt{par}$ operano
su porzioni disgiunte dello store, e lì la separazione non è una copia di
comodo.

Resta una sola eccezione, e non è una scorciatoia dell'esecutore ma l'ipotesi
sotto cui $\Inv{\mathtt{par}}$ è definita: quando un canale viaggia dentro un
messaggio i suoi titolari non sono più fissati dal testo, e il corpo inverso di
un $\mathtt{par}$ non si calcola in compilazione. Sono due programmi su
ottantanove, e per quei due l'inversione resta affidata al motore dinamico, che
accoppia gli appuntamenti mentre esegue invece di prevederli.

\paragraph{Il confronto che manca.} Il criterio di correttezza adottato qui è la
reversibilità con la confluenza che ne discende; la letteratura sui linguaggi
reversibili concorrenti usa invece \emph{causal safety} e \emph{causal
liveness}. Le due formulazioni parlano dello stesso fenomeno da posizioni
diverse, e metterle a confronto in modo puntuale, mostrando quale implichi quale
e sotto quali ipotesi, chiarirebbe dove si colloca la scelta di far comunicare i
rami soltanto per canali su porzioni disgiunte dello store.

\bigskip
\noindent
Reversibilità e concorrenza si studiano di solito separate: la prima nei
linguaggi di programmazione, la seconda nei calcoli di processo, con poche
occasioni di incontro. Kairos è un tentativo di averle insieme in un linguaggio
in cui si scrivono programmi, e il caso di studio è la verifica che
l'accostamento produca qualcosa: un unico programma che comprime e decomprime,
su blocchi trasformati in parallelo, è una cosa che nessuno dei due filoni
ottiene da solo.
